% Two Sources of Hardness in Generalized Heroes of Might and Magic III Combat
%
% This file mirrors paper/main.md, which is the master.  Edit main.md first.
%
% Compiles with tectonic (and with pdflatex; overfull-box warnings only).
% The long proofs are written out in full here: Theorem 3 in Appendix D, the
% lemma apparatus behind Theorems 1, 2, 4 and Proposition 1.1 in Appendix E.
% The proof repository proofs/ (in the artifact) retains the working documents.
%
% For a FUN / LIPIcs submission, replace the preamble below by
%
%   \documentclass[a4paper,USenglish,cleveref,autoref,thm-restate]{lipics-v2021}
%   \bibliographystyle{plainurl}
%
% and replace the hardcoded \newtheorem* environments by the numbered ones that
% lipics-v2021 defines (the labels below are hardcoded so that they literally
% match main.md: Theorem 1, Proposition 1.1, Theorem 2, Theorem 3,
% Corollary 3.1, Lemma 3.2, Theorem 4, Corollary 4.1, Corollary 4.2,
% Lemma 2.1, and Lemmas D.1-D.9 in Appendix D), and replace \maketitle/\author
% by \author{...}{...}{...}{...}{...}.

\documentclass[11pt,a4paper]{article}

\usepackage[T1]{fontenc}
\usepackage[utf8]{inputenc}
\usepackage[margin=3cm]{geometry}
\usepackage{amsmath,amssymb,amsthm}
\usepackage{booktabs}
\usepackage{array}
\usepackage{longtable}
\usepackage{alltt}
\usepackage{microtype}
\usepackage{url}
\usepackage[hidelinks]{hyperref}

% Latin Modern Roman has no italic small caps, so the \textsc problem names
% inside italic amsthm bodies fall back with ten font warnings; substitute the
% upright small-caps shape deliberately instead.
\DeclareFontShape{T1}{lmr}{m}{scit}{<->ssub*lmr/m/sc}{}

% The statement numbers below are hardcoded rather than counter-driven, because
% main.md numbers satellites off the theorem they hang from (Proposition 1.1
% under Theorem 1, Corollary 3.1 and Lemma 3.2 under Theorem 3, Corollaries 4.1
% and 4.2 under Theorem 4) while Lemma 2.1 is numbered by section.  No single
% amsthm counter produces that sequence, and the priority here is that the
% rendered labels match the master exactly.
\theoremstyle{plain}
\newtheorem*{thmone}{Theorem 1}
\newtheorem*{propdp}{Proposition 1.1}
\newtheorem*{thmtwo}{Theorem 2}
\newtheorem*{thmthree}{Theorem 3}
\newtheorem*{corthree}{Corollary 3.1}
\newtheorem*{lemresource}{Lemma 3.2}
\newtheorem*{thmfour}{Theorem 4}
\newtheorem*{corfourone}{Corollary 4.1}
\newtheorem*{corfourtwo}{Corollary 4.2}
\newtheorem*{lemnp}{Lemma 2.1}
% Appendix D runs its own D.1--D.9 sequence, hardcoded for the same reason.
\newtheorem*{lemreach}{Lemma D.1}
\newtheorem*{lemalt}{Lemma D.2}
\newtheorem*{lemadapt}{Lemma D.3}
\newtheorem*{lemembed}{Lemma D.4}
\newtheorem*{lemnoint}{Lemma D.5}
\newtheorem*{lembudget}{Lemma D.6}
\newtheorem*{lemconf}{Lemma D.7}
\newtheorem*{lemyes}{Lemma D.8}
\newtheorem*{lemcover}{Lemma D.9}
% Appendix E runs its own E.1--E.22 sequence, hardcoded for the same reason.
% E.8 is a Definition, so the plain sequence deliberately skips from E.7 to E.9.
\newtheorem*{lemEblow}{Lemma E.1}
\newtheorem*{lemEretal}{Lemma E.2}
\newtheorem*{lemEmetric}{Lemma E.3}
\newtheorem*{lemEsepone}{Lemma E.4}
\newtheorem*{lemEkillone}{Lemma E.5}
\newtheorem*{lemEyesone}{Lemma E.6}
\newtheorem*{lemEnoone}{Lemma E.7}
\newtheorem*{lemEtest}{Lemma E.9}
\newtheorem*{lemEfam}{Lemma E.10}
\newtheorem*{lemEseptwo}{Lemma E.11}
\newtheorem*{lemEacctwo}{Lemma E.12}
\newtheorem*{lemEyestwo}{Lemma E.13}
\newtheorem*{lemEnotwo}{Lemma E.14}
\newtheorem*{lemEseatsfour}{Lemma E.15}
\newtheorem*{lemErow}{Lemma E.16}
\newtheorem*{lemEroute}{Lemma E.17}
\newtheorem*{lemEreachfour}{Lemma E.18}
\newtheorem*{lemEreal}{Lemma E.19}
\newtheorem*{lemEaccfour}{Lemma E.20}
\newtheorem*{lemEyesfour}{Lemma E.21}
\newtheorem*{lemEnofour}{Lemma E.22}
\theoremstyle{definition}
\newtheorem*{defEmatch}{Definition E.8}
\theoremstyle{plain}
\theoremstyle{definition}
\newtheorem*{problem}{Problem}
\theoremstyle{remark}
\newtheorem*{remark}{Remark}

\newcommand{\ALLOC}{\textsc{Army-Allocation}}
\newcommand{\PLAY}{\textsc{Battle-Play}}
\newcommand{\PARTITION}{\textsc{Partition}}
\newcommand{\THREEPART}{\textsc{3-Partition}}
\newcommand{\PX}{\textsc{Planar-X3C}}
\newcommand{\SUBSETSUM}{\textsc{Subset-Sum}}
\newcommand{\src}[1]{\texttt{#1}}
\newcommand{\qedbox}{\hfill$\blacksquare$}

\title{Two Sources of Hardness in \\ Generalized Heroes of Might and Magic III Combat}

\author{Ivan Parfenchuk \\ \texttt{ivanparfenchuk@gmail.com}}
\date{}

\begin{document}
\maketitle

\begin{abstract}
We study the complexity of planning a single round of combat in a generalized version of
\emph{Heroes of Might and Magic III} (HoMM3)---the player's stacks start on prescribed deployment
hexes and the defending garrison follows a fixed scripted policy (wait, then defend)---using
the open-source VCMI reimplementation~\cite{VCMI26} as the specification of the rules. We show that one-round combat planning has \textbf{two
independent sources of strong NP-completeness}: pre-battle allocation of an army to slots is
strongly hard once the roster carries heterogeneous damage values, and in-battle target
selection remains strongly hard even with a fixed, single-type army. Between them we place a
positive result---a pseudo-polynomial $O(kB)$ dynamic program---which points at what makes
the tractable case tractable, and it is not the poverty of the roster but the triviality of
the \emph{reach structure} relating slots to enemies. All four hardness statements---three of
them strong, one weak and provably tight---hold under severe restrictions of the game: one round, flat damage, no spells, no heroes, no abilities,
and in three of the four no obstacles at all.

Every rule of the model is cited to a line of the VCMI source; the combat arithmetic and the
health mechanics are additionally cross-checked against the shipped engine classes
themselves. Each reduction was exhaustively regression-tested on bounded instances against an
executable transcription of the rules. That process caught three substantive errors in the
reductions and their verifiers, two of them in earlier versions of the proofs presented here;
external review caught others that bounded checks cannot see. We report them rather than
silently correcting them.
\end{abstract}

\begingroup\small
\noindent\textbf{MSC 2020.} 68Q17, 91A46.\\
\noindent\textbf{ACM CCS.} Theory of computation $\rightarrow$ Problems, reductions and
completeness; Applied computing $\rightarrow$ Computer games.\\
\noindent\textbf{Keywords.} computational complexity; NP-hardness; game AI; Heroes of Might
and Magic.
\endgroup

\section{Introduction}\label{sec:intro}

HoMM3 (1999) is a turn-based strategy game whose battles are fought on a hex grid between
armies of at most seven \emph{stacks}, each stack being a number of identical creatures.
Before a battle the player distributes creatures among the stacks; during a battle each stack
moves and attacks once per round. Both decisions are the ordinary business of playing the
game, and both are performed by hand, by every player, before and during every fight.

This paper asks how hard they are.

\subsection{The headline}\label{sec:headline}

\begin{quote}
\textbf{One-round combat planning in generalized HoMM3 has two independent sources of strong
NP-completeness: pre-battle allocation is strongly hard with heterogeneous damage types,
while in-battle target selection remains strongly hard even with a fixed, single-type army.}
\end{quote}

The two are independent in a precise sense. The first survives when the board wires each
enemy to exactly three deployment cells and the player merely fills the seats; the second
survives when the board wires nothing at all, and also when the board wires a great deal but
the army is a single creature type deployed one creature per slot. Neither theorem is a
restriction or special case of the other: each retains its hardness after the decision on
which the other construction rests has been fixed. A natural-looking conjecture that would
have unified them---that hardness scales with the diversity of the roster---is false; we state
it, and refute it, in Section~\ref{sec:thm3}.

Between the two we prove a positive result. In the family produced by our first reduction,
the problem is solvable by a dynamic program in $O(kB)$ arithmetic operations after a
polynomial preprocessing pass---pseudo-polynomial time---so on that family the weak hardness we
prove is exactly the hardness there is. Comparing that family to the hard ones isolates the
responsible parameter: the \textbf{reach hypergraph}, whose vertices are enemies and whose
hyperedges record which slots can engage which enemy. When it is a perfect matching---and
every enemy stack is a single creature, so that a blow either finishes its target or scores
nothing---the allocation separates into independent per-slot threshold decisions and the
problem is a knapsack (Proposition 1.1 states the exact hypotheses; with multi-creature
stacks partial kills score and the threshold framing fails). Every hard family we construct
breaks the matching---a slot engaging several
enemies, or several slots the same enemy---and strong hardness then appears with a rich
roster, with a poor one, and with no allocation decision at all.

We do not claim a dichotomy. We have hard cases on both sides of the roster axis and a
tractable case with a degenerate reach structure; that is a description of the examples we
have, not a theorem about all restrictions---trivially easy non-matching reach structures
exist (an instance consisting of a single enemy reached by two slots, say), so ``as soon as the matching breaks, hardness
appears'' would be false and we do not claim it.

\subsection{Contributions}\label{sec:contrib}

\begin{enumerate}
\item \textbf{A formal model of HoMM3 combat cited line by line to the VCMI engine}
  (Section~\ref{sec:model}), with an executable transcription, and with the combat arithmetic
  and health mechanics cross-checked against the shipped engine classes.
\item \textbf{Theorem 1} (Section~\ref{sec:thm1}): allocation is NP-complete with a
  \emph{single} creature type, $R=1$, no obstacles, one creature per enemy stack, and a
  one-row battlefield. Weak hardness, from \PARTITION~\cite[SP12]{GJ79}.
\item \textbf{Proposition 1.1} (Section~\ref{sec:thm1}): on the family the reduction
  \emph{constructs}---which additionally has persistent matching reach; it and Theorem 1's
  one-creature-per-enemy-stack restriction are both load-bearing---an $O(kB)$ dynamic program solves the problem, so
  Theorem 1 is tight and strong hardness there is impossible unless
  $\mathrm{P}=\mathrm{NP}$.
\item \textbf{Theorem 2} (Section~\ref{sec:thm2}): allocation is \emph{strongly} NP-hard, from
  \THREEPART---NP-complete in the strong sense~\cite[SP15]{GJ79}---with every \emph{player}
  creature type of stock one and one creature per enemy stack.
\item \textbf{Theorem 3 and Corollary 3.1} (Section~\ref{sec:thm3}): the problem is strongly
  NP-hard with a \textbf{single creature type}, and remains so with the allocation fixed to
  one creature per slot. The reduction is from planar exact cover by 3-sets and
  uses the engine's lower clamp on damage as its source of arithmetic. This refutes the
  roster-diversity conjecture.
\item \textbf{Theorem 4 with Corollaries 4.1 and 4.2} (Section~\ref{sec:thm4}): strong
  NP-hardness survives three further restrictions---an obstacle-free rectangle with
  \emph{complete reachability} (in the starting position, every stack can attack every enemy);
  a given allocation, which isolates unrestricted target assignment as the hard decision; and
  an objective changed from creatures killed to hit points removed.
\item \textbf{A machine-checking methodology} (Section~\ref{sec:checking}) and an honest
  report of the errors it caught. One scope note belongs up front: the checks exercise the
  constructions on bounded instances; the embedding algorithm of Lemma D.4 is implemented
  step by step---with three named subroutine substitutions, disclosed in
  Appendix~\ref{app:mcscope}---and
  validated on the whole instance corpus (Appendix~\ref{app:mcscope}; the suite table is
  the artifact's \src{README.md}), but its
  correctness on \emph{all} inputs---like every proof in this paper---rests on the hand
  proof, not on a proof assistant.
\item \textbf{An empirical companion note} in the artifact
  (\src{paper/companion-empirics.md}) measures the certified optimum against one-shot
  allocations proposed by three tiers of language models on 145 instances, each allocation
  completed by oracle-optimal play; it is not part of this paper's claims.
\end{enumerate}

\subsection{What is not new}

The combinatorial cores are standard. Theorem 2 and Theorem 4 use the \THREEPART{} targeting
structure already seen in Hearthstone puzzles~\cite{HLW20}; Theorem 3 is planar exact-cover
targeting; the homogeneous-resource split of Theorem 1 is the mechanism of
\cite[Thm.~7]{FB10}, published in 2010. We claim no new reduction technique. We claim
three things: a \emph{separation}---contrasting restricted families that locate where HoMM3's
hardness resides; the fact that the decisions in question are literal player-facing actions of
a shipped game rather than modelling devices; and a standard of verification uncommon in this
genre. Section~\ref{sec:related} sets out the neighbours in detail.

\section{The model}\label{sec:model}

\subsection{Why the game must be generalized}\label{sec:generalized}

Shipped HoMM3 has a fixed battlefield of $11\times17$ hexes, at most 7 army slots per side
(rule R14, Appendix~\ref{app:rules}), and a fixed creature roster. A game with a bounded state
space is decidable in constant time, so we follow the standard practice of the genre and
generalize: the battlefield is $n\times m$, the number of slots is $k$, and creature
statistics are arbitrary integers given in the input. Numbers are written in binary unless
stated otherwise; binary is the honest encoding for a game that routinely carries stack sizes
in the tens of thousands.

One further thing is generalized, and it deserves its own flag. \textbf{The deployment cells
are part of the input.} In the shipped game a stack's starting hex is a function of the slot
count and slot index, read from fixed formation tables (R14); in our problem the scenario
prescribes them, as it prescribes obstacles and enemy placement. The object we study is
therefore \emph{generalized battle-scenario planning with prescribed deployment cells}, and a
claim that ``only the bounds are generalized'' would be false. All four theorems use this
freedom; Theorem 3 uses it heavily. Recovering any of them under native formations is open
(Section~\ref{sec:open}).

A second, smaller generalization was found by external review, and we record it rather than
let it be found again. The engine does \textbf{not} make every cell of its rectangle usable: the first and
last columns fail \src{BattleHex::isAvailable()} and are labelled \src{SIDE\_COLUMN} by
\src{getAccessibility} (R16). Our constructions do use column 0---Theorem 1 puts $p_1$ there,
Theorem 2 deploys its first group in it ($q_1^1=(0,1)$ is a deployment hex, which under R16
could not be occupied at all), and Theorem 4 puts $p_1$ in it (its line
$p_j = (j-1, 0)$ runs along the top row, so only $p_1$ touches column 0). So ``the battlefield is
$n\times m$'' is generalized in one further respect: we treat every cell as usable. The repair is
mechanical: pad with one unusable column on each side and shift every constructed coordinate
one \textbf{column} to the right---the row must not change, since shifting $y$ flips row parity
and is not adjacency-preserving on the offset grid.
\textbf{Speeds must be left alone---for Proposition 1.1.} An earlier version of this
paragraph also prescribed widening the speed bounds by two. Translation preserves every hex
distance, so both directions of Theorem 1 carry over unchanged (its no-direction is in fact
reach-independent; Appendix~\ref{app:thm124}); but raising the player's speed from 2 to 4
raises its strike radius from 3 to 5, and on the \PARTITION{} instance $(2,4)$ the second
slot then reaches $E_1$ as well as $E_2$, so the family no longer has the persistent matching
reach that Proposition 1.1 assumes. (We checked this against the executable model rather than
by inspection.) For Theorem 3, horizontal translation preserves row parity, adjacency, distances and
components, so (I1)--(I3) survive and (I4) does too since no region grows; if $\sigma$ is defined
as the board's hex count it grows by twice the height, not by two. The shift has not been carried
out in the constructions or the verifiers as they stand, so we flag it here as a gap between the
model and the engine rather than claim it away.

\subsection{The rules, with citations}\label{sec:rules}

We use VCMI~\cite{VCMI26}, an open-source reimplementation whose combat module reproduces the original
game's damage numbers, as the specification. Every rule below carries a tag \texttt{R1}--\texttt{R17};
Appendix~\ref{app:rules} maps each tag to the exact \src{file:line} of the checkout (commit
\src{b5cee70}), so the body can be read without recourse to the paths. The full model is
\src{MODEL.md}; this is the fragment the theorems need.

\paragraph{Battlefield.} A hex grid in offset (``even-row shifted'') coordinates with six
neighbour directions (R1); distance is computed through axial coordinates (R2). A subset of
hexes may be impassable, merging the engine's three mechanisms---battlefield
\src{impassableHexes}, obstacle objects, and siege walls (R3). \textbf{A hex occupied by a
living unit is not enterable, and a dead unit stops blocking} (R4), which matters a great deal
in Section~\ref{sec:thm3}.

\paragraph{Units.} A creature type is
$(\mathit{att},\mathit{def},\mathit{dmg}_{\min},\mathit{dmg}_{\max},\mathit{hp},
\mathit{spd},\mathit{flags})$, plus a nonnegative \emph{value} used by the objective. A stack
is a type together with a creature count, and its health is stored as a pair
$(\mathit{fullUnits},\mathit{firstHPleft})$ representing a pool
$\mathit{avail}=\mathit{firstHPleft}+\mathit{hp}\cdot\mathit{fullUnits}$ (R5); a fresh stack of
$c$ creatures has $\mathit{fullUnits}=c-1$ and $\mathit{firstHPleft}=\mathit{hp}$, so
$\mathit{avail}=c\cdot\mathit{hp}$ and $\mathrm{count}=c$. The rule the
reductions live on is the \textbf{effective count} (R6):
\[
  \mathrm{count}(S)=\mathit{fullUnits}+[\mathit{firstHPleft}>0].
\]
A stack's offensive output is proportional to $\mathrm{count}$, not to $\mathit{avail}$. A
stack of one Archangel reduced to 1 hit point out of 250 deals exactly as much damage as a
healthy one. Damage output is a \textbf{step function} of damage taken.

\paragraph{Damage (R7).} With $\Delta=\mathit{att}(\text{attacker})-\mathit{def}(\text{defender})$:
\begin{align*}
  f_{\mathrm{att}} &= 1+\min(0.05\Delta,\,4.0) \quad\text{if } \Delta>0,\ \text{else } 1\\
  f_{\mathrm{def}} &= 1-\min(0.025(-\Delta),\,0.7) \quad\text{if } \Delta<0,\ \text{else } 1\\
  \mathrm{dmg} &= \max\bigl(1,\ \lfloor \mathrm{count}\cdot d\cdot f_{\mathrm{att}}\cdot
      f_{\mathrm{def}}\rfloor\bigr)
\end{align*}
In the mathematical model the constants are the exact rationals $1/20$, $1/40$, $4$, $7/10$
and the floor is exact integer arithmetic; the engine evaluates the same formula in IEEE-754
doubles, and Section~\ref{sec:floats} documents the one divergence the cross-check has
found---a claim scoped to the cases run, not a proof that no other exists. The reference
simulator also computes in floats, and floats and exact rationals do not agree on every call
either (Section~\ref{sec:floats}); the theorems consume only what holds under both, and the
certificates of the empirical companion note are re-derived under exact arithmetic
(Section~\ref{sec:floats}). Two
features of this formula do real work below. The multiplicative dependence on
$\mathrm{count}$ makes damage linear in stack size; the \textbf{lower clamp at 1} breaks that
linearity from below, so that splitting a stack into singletons can deliver strictly more
total damage than keeping it whole. Section~\ref{sec:thm3} is built entirely on the clamp.

\paragraph{Kills and overkill (R8).} Damage accumulates in the defender's pool; excess beyond
the pool is discarded (the \emph{overkill rule}), and
\[
  \mathrm{kills}(D)=\begin{cases}
    0 & \text{if } D<\mathit{firstHPleft}\\[2pt]
    \min\!\left(1+\left\lfloor\dfrac{D-\mathit{firstHPleft}}{\mathit{hp}}\right\rfloor,
      \mathrm{count}\right) & \text{otherwise.}
  \end{cases}
\]

\paragraph{Turn structure.} Play proceeds in rounds; within a round each living stack acts once in decreasing
order of $\mathit{spd}$, ties broken by side then slot index (R9). That tie rule is a
\textbf{simplification}: the engine alternates sides on equal initiative according to which side
moved last, and gives the attacker priority only on the first turn (R17). We use the simplified
rule because in every construction here each player stack is strictly faster than every enemy, so
no equal-initiative tie between the sides ever arises and the two rules coincide on the instances
we build; but the simplification is real and the executable transcription shares it, so the
mechanics tests cannot detect the difference. What the tests cannot detect, the engine harness
pins directly: two dedicated tie cases run the engine's own queue and confirm the alternation
rule is exactly as stated here (Section~\ref{sec:layers}). \textbf{Speed is simultaneously
initiative and movement range} (R10)---a construction cannot set them independently, which rules
out the obvious ``slow but far-reaching'' gadget. Stacks that use \texttt{WAIT} act after all
\texttt{NORMAL} stacks, in \emph{increasing} speed order (R9)---the fact the garrison policy of
Section~\ref{sec:policy} rests on, and one the engine harness executes through the engine's own
queue (Section~\ref{sec:layers}). Movement is unweighted BFS over enterable hexes bounded by
$\mathit{spd}$; \texttt{WALK\_AND\_ATTACK} moves to a hex adjacent to the target and strikes
(R11), so a melee stack of speed $s$ can strike only enemies at distance at most $s+1$;
whether it can actually strike such an enemy depends on the free hexes available.

\paragraph{Retaliation (R12).} The attacker's blow resolves first; the defender retaliates
afterwards if it is still alive, has a charge left, and the attacker lacks
\texttt{BLOCKS\_RETALIATION}. \textbf{A stack killed outright does not retaliate.} Charges
reset at the round boundary, so the first attacker into a stack absorbs the retaliation and
later attackers in the same round strike free.

\paragraph{Determinization.} We restrict to instances with
$\mathit{dmg}_{\min}=\mathit{dmg}_{\max}$, exclude morale, luck, spells, heroes, and every
conditional damage multiplier, and call the result \texttt{H3-det} (\src{MODEL.md} \S7).
Constraining the \emph{instances} rather than reinterpreting the \emph{rules} is deliberate:
\texttt{H3-det} is then a genuine special case of the generalized model of
Section~\ref{sec:generalized}---not of the shipped game, whose board, roster and formations are
fixed---so hardness transfers upward to that model, and no rule was reinterpreted to make a
proof work.
One further restriction deserves its own name. Every construction in this paper uses only
melee, single-hex creatures with the default single retaliation charge---the fragment
\texttt{H3-det-melee} (\src{MODEL.md} Definition~7.1a), on which the damage formula above
is the \emph{whole} formula: no ranged or distance penalties, no breath or multi-target
attacks, no double-wide movement cases. Our hardness results are therefore statements about
\texttt{H3-det-melee}, and they transfer upward to \texttt{H3-det} and to the generalized
model because a restriction of the instance family only strengthens a hardness claim; the
instances of the empirical companion note (Section~\ref{sec:contrib}) live in the same fragment.

\subsection{The problem}\label{sec:problem}

\begin{problem}[\ALLOC]
\textbf{Input.} A battlefield $(n,m,\mathit{obstacles})$; $k$ slots with deployment hexes
$p_1,\dots,p_k$; a multiset $A$ of player creatures as (type, count) pairs; a fixed
defence---enemy stacks with types, counts and hexes---which plays the fixed policy
$(\ddagger)$ of Section~\ref{sec:policy}; a round bound $R$ in unary; a target
$W\in\mathbb{Z}_{>0}$. Every creature type carries a nonnegative integer \emph{value} as part
of its tuple (Section~\ref{sec:rules}); the objective reads the values of enemy types only.
The deployment hexes $p_1,\dots,p_k$ are pairwise distinct, passable, and distinct from every
enemy hex, and an allocated stack begins the battle on the deployment hex of its slot. Where
R9's speed comparison ties across the two sides, the player's stack acts first.

\noindent\textbf{Question.} Is there an allocation of $A$ to the $k$ slots (each slot receiving
at most one type, each type's total at most its stock) and a sequence of player actions such
that after $R$ rounds against $(\ddagger)$ the total value of enemy creatures killed is at
least $W$?
\end{problem}

\PLAY{} is the same question with the allocation given as part of the input, so \ALLOC{}
contains \PLAY{} as a special case.

Three points on the formulation. The enemy plays one concrete scripted policy, not an
adversarial one, which keeps the problem in NP and matches the informal ``against a fixed
defence''; the adversarial version is a different problem. Baking $(\ddagger)$ into the problem
rather than taking an arbitrary ``deterministic poly-time policy $\pi$'' as input is
deliberate: an arbitrary encoded program with a promised running time is not a syntactically
checkable input restriction, a circuit encoding would repair that at no gain, and every theorem
uses $(\ddagger)$ anyway---hardness for the one fixed policy is the stronger statement. The
objective counts \textbf{whole creatures killed}, weighted by the per-type value. That is the
game's own accounting: a stack at one hit point fights at full strength, so hit points removed
without a kill buy the player nothing in that round. It is \emph{not}, however, the sole source
of hardness---counting hit points removed collapses only the matching-reach family of Theorem 1
to a separable concave sum; on the family of Theorem 4 the hit-point objective remains strongly
NP-hard (Corollary 4.2). And $R$ is given in unary so that evaluating a certificate is
polynomial.

\begin{lemnp}[membership]
$\ALLOC\in\mathrm{NP}$.
\end{lemnp}

\noindent\emph{Proof.} Fix the encoding: the battlefield is listed hex by hex and $R$ is unary.
A certificate is the allocation together with the player's actions: for each stack-round pair,
an optional \texttt{WAIT} bit and one terminal action: \texttt{DEFEND}, a destination hex
(move), or a (destination hex, target) pair (\texttt{WALK\_AND\_ATTACK})---at most $2kR$
action tokens, since a stack that waits still takes its terminal action later in the same round
(Section~\ref{sec:rules}). Simulating the battle takes a polynomial number of arithmetic
operations and BFS computations, $(\ddagger)$ is computable in constant time, and comparing
destroyed value to $W$ completes the check. \qedbox

\subsection{The garrison policy}\label{sec:policy}

All four theorems use the same scripted defence, and it has to be pinned down precisely,
because ``hold position'' does not determine an action: \texttt{H3-det} retains movement,
\texttt{WAIT}, and \texttt{DEFEND}.

\begin{quote}
$(\ddagger)$ If the stack has not waited this round, issue \texttt{WAIT}; on its postponed
activation, issue \texttt{DEFEND} at its current hex.
\end{quote}

Both are shipped actions. \texttt{WAIT} postpones the stack's terminal action into the round's
\texttt{WAIT} phase, which runs after all \texttt{NORMAL}-phase activations, in \emph{increasing}
speed order (R9); \texttt{DEFEND} (R13) ends the turn without moving or attacking and grants
$+20\,\%$ defence---an integer bonus with a floor of $+1$---until the stack next receives a turn.
Three consequences are used throughout:

\begin{itemize}
\item no enemy ever initiates an attack, so no player stack ever delivers retaliation damage;
\item a waiting or defending enemy still retaliates when struck---neither action consumes the
  retaliation charge;
\item \textbf{the one-round lemma}: in round 1, the blow of a player stack that did not wait
  lands at its \texttt{NORMAL}-phase activation, before any enemy's postponed \texttt{DEFEND},
  and therefore meets the un-raised defence---\emph{regardless of relative speeds}; and every
  blow, waiting or not, delivers at most its nominal damage, since raising the target's defence
  can only lower $f_{\mathrm{def}}$. \textbf{Every statement quantifying over an arbitrary play
  is therefore phrased as ``at most the nominal damage''}; equalities are asserted only for the
  constructed witness plays, which never wait.
\end{itemize}

The lemma's scope is exactly one round, and both ways it fails outside that scope are worth
stating, because an earlier justification (``no blow ever meets the bonus'') claimed too much and
was refuted by external review. A player stack that itself waits is scheduled in the
\texttt{WAIT} phase by \emph{increasing} speed, so a slower enemy's postponed \texttt{DEFEND} can
land first and the player's postponed blow then meets the bonus---for such blows only the
inequality survives. And for $R\ge 2$ the bonus persists past the round boundary until the
enemy's next activation, so a fast player striking early in round 2 meets the bonus left over
from round 1. All four theorems set $R=1$; a multi-round extension must redo this analysis.

An earlier version of $(\ddagger)$ issued \texttt{DEFEND} at the stack's own turn. For the
theorems the two are interchangeable---every constructed player stack is strictly faster than
every enemy---but they part company on the corpus of the empirical companion note, whose
natural instances contain enemies \emph{faster} than the player, and there the old policy made
six recorded optima unattainable (the companion note gives the history and the machine
checks that close it). The present $(\ddagger)$ was adopted after being verified reduction by
reduction by external review; the mechanics it stands on---the \texttt{WAIT}-phase order, the
bonus arithmetic and duration, and the retaliation-charge neutrality of both actions---are
unit-checked against the cited engine lines, and the searches of Section~\ref{sec:checking}
include a variant in which the defence executes $(\ddagger)$ literally, phase by phase
(Section~\ref{sec:layers}).

That the policy is load-bearing is not obvious. The first version of Theorem 1 used an
attacking garrison and was wrong: an enemy that attacks provokes a \emph{retaliation}, which
delivers a second blow to that same enemy inside the same round, and the reduction then decides
\SUBSETSUM{} with the wrong budget. Section~\ref{sec:errors} records it.

\section{Results}\label{sec:results}

We present the theorems in pedagogical order. The two-source claim of
Section~\ref{sec:headline} is assembled from Theorems 2, 3 and 4 at the end of
Section~\ref{sec:thm4}.

Throughout, $(\star)$ denotes the specialization of the damage formula used by Theorems 1, 2
and 4: every player type and every enemy type has
$\mathit{att}=\mathit{def}=\alpha$ with $\alpha:=1$, and every enemy type has flat damage 1.
Then $\Delta=0$ in both directions, $f_{\mathrm{att}}=f_{\mathrm{def}}=1$, and a stack of $c$
creatures of flat per-creature damage $d$ delivers nominal damage $c\,d$. Retaliation costs the
player nothing under $(\star)$: the defence never initiates (Section~\ref{sec:policy}), so a
player stack takes at most one retaliation, and it takes it \emph{after} its own blow (R12),
so no blow already delivered is affected. Theorem 3 deliberately breaks $(\star)$ and uses the
defence factor.

\subsection{Theorem 1, and an algorithm that matches it}\label{sec:thm1}

\begin{thmone}
\ALLOC{} is NP-complete, already for instances with $R=1$, a \textbf{single creature type} in
the player's army, one creature per enemy stack, no obstacles, and a battlefield of one row.
\end{thmone}

The interesting word is \emph{single}. There is nothing to choose about what to bring; the only
free variable is how many creatures go in each slot.

\smallskip
\noindent\emph{Construction.} From a \PARTITION{} instance~\cite[SP12]{GJ79}
$a_1,\dots,a_n$ with $\sum a_i=2B$: one row of $5n$ hexes; block $j$ occupies hexes
$5(j-1),\dots,5j-1$, with the deployment hex $p_j:=5(j-1)$ and the enemy $E_j$ at
$e_j:=5(j-1)+1$. Under $(\star)$ the player has one type with
$\mathit{att}=\mathit{def}=1$, flat damage 1, $\mathit{hp}=5$, speed 2, no flags, value $0$
and stock exactly $B$; $E_j$ is one creature of a type with $\mathit{att}=\mathit{def}=1$,
flat damage 1, $\mathit{hp}=\mathit{value}=a_j$, speed 1 and no flags. Set $R=1$, $W=B$.

\smallskip
\noindent\emph{Correctness.} A stack of speed $s$ can strike only enemies at hex distance at
most $s+1$, in any position and regardless of which hexes are free (Lemma~E.3); consecutive
deployment hexes are 5 apart, so every foreign enemy sits at distance at least 4 from $p_j$
against a strike radius of $\mathit{spd}+1=3$, and slot $j$ can strike $E_j$ and nothing
else---in every position of the play, not only at the start, because the obstruction is
distance rather than blocking, and deaths do not change distances. A stack strikes at most
once, so the striker sets of distinct enemies are disjoint, and the total damage $E_j$
absorbs is at most $c_j$ (Section~\ref{sec:policy}); $E_j$ dies only if $c_j\ge a_j$, and a
non-waiting blow with $c_j\ge a_j$ kills. If \PARTITION{} has a solution $S$, allocate
$c_j=a_j$ for $j\in S$: total $B$, destroyed value $B$. Conversely, if the game is a
yes-instance with dead set $S$, then
$B\le\sum_{j\in S}a_j\le\sum_{j\in S}c_j\le\sum_j c_j\le B$, so every inequality is tight and
$S$ solves \PARTITION. On a malformed encoding, a non-positive $a_i$, or an odd $\sum a_i$, the
reduction outputs the fixed no-instance of Lemma D.4; on the empty instance $n=0$---a
\PARTITION{} \emph{yes}-instance---it outputs $G((1,1))$, which keeps the map total. Theorems~2 and~4 route $|a|\ne 3m$, some $a_i\le 0$, $\sum a_i\ne mT$ or an $a_i$
outside $(T/4,T/2)$ to $G_{\mathrm{no}}$ the same way (Appendix~\ref{app:thm124}, the
Theorem 2 and Theorem 4 sections); $G_{\mathrm{no}}$ is a $1\times 1$ board and not of the
shape Theorem 4 names---its family clause describes the image of every well-formed encoding,
and the degenerate encodings, no-instances certified in polynomial time, go to the one fixed
no-instance shared by all three constructions. With Lemma 2.1 this gives NP-completeness. The full proof is Appendix~\ref{app:thm124}. \qedbox

\smallskip
The numbers $a_j$ are carried as hit points in binary, so this is only \emph{weak}
hardness---and that is exactly right, because the family admits an algorithm that matches
that bound:

\begin{propdp}
On the family of Theorem 1---single creature type of flat damage $d\ge1$, $R=1$, policy
$(\ddagger)$, damage under $(\star)$, \textbf{one creature per enemy stack}, and
\emph{persistent matching reach}: a bijection $j\mapsto E_j$ such that, whatever the
feasible allocation, in every position of
round 1 reachable by legal play from its starting position, for every slot $j$ whose stack has \emph{not yet taken its
terminal action}, the set of enemies that stack can strike is exactly
$\{E_j\}$---\ALLOC{} is decidable in $O(k\,B)$ arithmetic operations and $O(B)$ working space
after a polynomial-time preprocessing pass (one breadth-first search per slot), hence in
pseudo-polynomial time; $B$ is the stock.
\end{propdp}

\noindent\emph{Proof.} After deployment the play is forced: slot $j$ either finishes $E_j$ or
achieves nothing. This is where \textbf{one creature per enemy stack} is load-bearing: against
a stack of several creatures a non-finishing blow still kills whole creatures and still
scores, the per-slot value is a staircase in $c_j$ rather than a threshold, and the 0-1
framing below is false (open problem 4, Section~\ref{sec:open}---a later review pass
exhibited a matching-reach instance with six-creature stacks where the threshold rule returns
$0$ and the true optimum is $3$). Against a single creature of $t_j$ hit points, slot $j$ scores $v_j$
iff its nominal damage $c_j\,d$---delivered exactly by a non-waiting blow under $(\star)$,
since $\Delta=0$ makes both factors $1$; a waiting blow meets the postponed \texttt{DEFEND}
bonus and delivers at most nominal, which only helps the upper-bound
direction---reaches $t_j$, i.e.\ iff $c_j\ge b_j:=\lceil t_j/d\rceil$ (well defined since
$d\ge1$), and surplus is wasted. The optimum is
$\max\{\sum_{j\in S}v_j:\sum_{j\in S}b_j\le B\}$, a 0-1 knapsack over $k$ items, solved by
the textbook dynamic program. The full proof, with the hypothesis consumed exactly where it
is needed, is Appendix~\ref{app:thm124}. \qedbox

\smallskip
Proposition 1.1 is machine-checked twice: the knapsack program against exhaustive subset
enumeration of the same abstraction on 2000 random instances, and---since round 10---against
exhaustive play of built corridor instances in the game model itself, including a negative
control with multi-creature enemy stacks on which the suite asserts that the threshold rule
and the game \emph{disagree} (\src{dp\_single\_type.py}). It matters for three reasons. It
shows Theorem 1 is tight. It disposes of an apparent contradiction---``Theorem 1 gives a DP,
Theorem 2 forbids one''---which confuses two incomparable restrictions of the same general
problem. And, read against Theorems 3 and 4, it identifies the parameters that are actually
responsible: not the roster, but the matching reach structure together with single-creature
stacks (open problem 4 asks what survives without the latter).

\subsection{Theorem 2: allocation-driven strong hardness}\label{sec:thm2}

\begin{thmtwo}
\ALLOC{} is \textbf{strongly} NP-hard, already for $R=1$, no obstacles, no abilities, one
creature per enemy stack, and instances in which every \emph{player} creature type has stock
exactly one.
\end{thmtwo}

\smallskip
\noindent\emph{Construction.} From a \THREEPART{} instance $(a_1,\dots,a_{3m};T)$ with
$\sum a_i=mT$ and $T/4<a_i<T/2$---NP-complete in the strong sense~\cite[SP15]{GJ79}---take
three rows and $8m+2$ columns. For each group $g$, an enemy $E_g$ with
$\mathit{att}=\mathit{def}=1$, flat damage 1, $\mathit{hp}=T$, speed 1, no flags and value 1
sits at $(X_g,1)$ with $X_g=8(g-1)+1$, and its three deployment hexes are three of the six
neighbours of $(X_g,1)$, named explicitly: row 1 is odd, and
$q_g^1=(X_g-1,1)$, $q_g^2=(X_g,0)$, $q_g^3=(X_g,2)$, all inside the three-row board. The
player has $3m$ types, type $C_i$ with $\mathit{att}=\mathit{def}=1$, flat damage $a_i$,
$\mathit{hp}=5$, speed 2, no flags, value 0 and \textbf{stock one}. $R=1$, $W=m$.

\smallskip
\noindent\emph{Correctness.} Each $q_g^r$ is adjacent to $E_g$ and at hex distance at least 7
from every other enemy (tight, attained by $q_{g+1}^1$ against $E_g$) against a strike radius
of $\mathit{spd}+1=3$, so the three slots of group $g$ strike $E_g$ and nothing else. Each
stack holds \emph{at most} one creature (stock one; slots may stay empty), a stack strikes at
most once, and under $(\ddagger)$ no enemy initiates, so no player stack ever delivers a
retaliation blow (Lemma E.2); hence the damage available to $E_g$ is at most
$\sum_{i\in S_g}a_i$ where $S_g$ is the set of types allocated to group $g$'s slots, and the
$S_g$ are pairwise disjoint. Killing all $m$ enemies forces $\sum_{i\in S_g}a_i\ge T$ for every
$g$; summing against $\sum a_i=mT$ makes every inequality tight, and $T/4<a_i<T/2$ forces
$|S_g|=3$. The sufficiency direction, where the game semantics are discharged---all three
stacks of a group strike in the \texttt{NORMAL} phase before the enemy's postponed
\texttt{DEFEND}, and their damage accumulates in the health pool (R8)---is Lemma E.13; the
degenerate encodings are routed as in Theorem 1. The full proof is
Appendix~\ref{app:thm124}. \qedbox

\begin{remark}
\textbf{The three rows are not decoration.} An earlier version of this construction used a
single row, as Theorem 1 does. That is impossible: in one row a hex has two neighbours, so a
third stack cannot reach $E_g$ without walking through a hex occupied by an ally, and occupied
hexes are not enterable. The bug was found by a geometry self-check that verifies the
reachability lemma on the \emph{built} instance with all slots occupied (the geometry
self-check of the Theorem 1--2 suite, Section~\ref{sec:layers}).
\end{remark}

\subsection{Theorem 3: a single type is enough}\label{sec:thm3}

Theorems 1 and 2 together tempt one into a story: hardness scales with the diversity of the
roster, since the single-type family of Theorem 1 has a pseudo-polynomial algorithm and strong
hardness needed $3m$ distinct types. An earlier draft of this work recorded that story as an
explicit conjecture. It is false.

\begin{thmthree}
\ALLOC{} is \textbf{strongly} NP-hard, already for $R=1$, a \textbf{single creature type} in
the player's army, one creature per enemy stack, all enemy creatures of one type and value 1,
flat damage, no abilities, and static impassable hexes as the only terrain feature. With
Lemma 2.1 it is strongly NP-complete on this family.
\end{thmthree}

\begin{corthree}[fixed allocation]
The same instances remain strongly NP-hard when the allocation is given: fix one creature in
every slot. \PLAY{} is strongly NP-hard on single-type instances.
\end{corthree}

\noindent\textbf{Corollary 3.1 is the form to quote, for a reason worth stating plainly.} In
every yes-instance of Theorem 3 the winning allocation is unique and is the all-ones
vector---the budget forces $3q$ creatures into $3q$ slots, one each. There is therefore no
interesting sizing decision, and it would be wrong to advertise this as ``single-type stack
sizing is strongly hard''. What the theorem establishes is that \textbf{target selection on the
reach hypergraph is strongly hard}, and that it stays hard when the roster is reduced to one
type and the allocation is removed from the problem altogether. That is precisely what kills
the roster-diversity conjecture.

\paragraph{Source problem.} \PX: exact cover by 3-sets whose element/set incidence graph is
planar. Dyer and Frieze prove exactly this problem NP-complete---their Lemma 2.2 states
``Planar X3C is NP-complete'', with planarity defined on precisely this incidence
graph~\cite[p.~175]{DF86}---and their instances additionally have every element in two or
three sets (p.~178), so the hardness survives that restriction too.

\paragraph{The arithmetic.} One player type $P$ with $\mathit{att}=1$, $\mathit{def}=1$, flat
damage 1, $\mathit{hp}=4$, speed $\sigma$ equal to the board's hex count, stock exactly $3q$;
one enemy type $Q$ with $\mathit{att}=1$, $\mathit{def}=\mathbf{27}$, flat damage 1,
$\mathit{hp}=3$,
speed 1, value 1. Then $\Delta=-26$ and $0.025\cdot 26=0.65$, below the cap $0.7$, so
$\mu:=0.35$ and a player stack of $c$ creatures delivers nominal damage
\[
  D(c)=\max(1,\lfloor\mu c\rfloor):\qquad
  D(1),\dots,D(12)=1,1,1,1,1,2,2,2,3,3,3,4.
\]

\begin{lemresource}[resource lemma]
Let $0<\mu<1$ and $D(c)=\max(1,\lfloor\mu c\rfloor)$. If stacks of sizes
$c_1,\dots,c_r\ge 1$ each deliver at most $D(c_i)$ to one target and the total is at least 3,
then $\sum_i c_i\ge 3$, with equality iff $r=3$ and $c_1=c_2=c_3=1$.
\end{lemresource}

\noindent\emph{Proof.} For $c=1$, $D(1)=1=c$; for $c\ge 2$, $\lfloor\mu c\rfloor\le\mu c<c$
and $1<c$, so $D(c)<c$. Hence $3\le\sum_i D(c_i)\le\sum_i c_i$. If $\sum_i c_i=3$ then
$\sum_i D(c_i)=\sum_i c_i$, so $D(c_i)=c_i$ and thus $c_i=1$ for every $i$; since
$\sum_i c_i=3$, $r=3$. \qedbox

\smallskip
This is the damage floor doing the work. A lone creature always delivers a full point no matter
how outclassed; $c$ creatures in one stack deliver $\lfloor\mu c\rfloor<c$. Splitting into
singletons is strictly better---a real if inelegant HoMM3 tactic---and the granularity it buys
is what a covering problem needs. Concretely, three damage costs 9 creatures in one stack, 7 in
two, and 3 in three.

Stating the lemma for all $\mu\in(0,1)$ rather than for $0.35$ is deliberate, because the
construction's two branches sit on opposite sides of the defence cap. An undefended blow sits
at $0.025\cdot 26=0.65$, clear of the cap. A defending enemy (Section~\ref{sec:policy}) has
defence 32, hence $0.025\cdot 31=0.775$---\emph{past} the cap, so that branch is clamped to
$0.7$ and $\mu=0.3$; and VCMI's hand-written JSON parser loads the cap constant $0.7$ as
$0.7000000000000001$ (Section~\ref{sec:floats}), so the clamped branch lives in two slightly
different arithmetics, engine and model disagreeing by one unit in the last place (ULP)
exactly at integer boundaries
of $\mathit{base}\times 0.3$. The lemma is indifferent to all of it: every $\mu$ in play---$0.35$,
$0.3$, or the engine's $0.29999999999999993$---lies in $(0,1)$, which is the only property the
proof consumes.

\paragraph{The board.} Four invariants carry the whole correctness argument, and nothing below
refers to the layout:

\begin{itemize}
\item \textbf{(I1)} every enemy hex $z_S$ has exactly three free neighbours, pairwise
  non-adjacent;
\item \textbf{(I2)} the free non-enemy hexes fall into exactly $3q$ connected components, one
  per element---the \textbf{region} $R_e$;
\item \textbf{(I3)} $R_e$ holds the deployment hex $p_e$ and exactly the dockings $d_S^e$ for
  $S\ni e$, each adjacent to $z_S$ alone;
\item \textbf{(I4)} the speed $\sigma$ satisfies $\sigma\ge\max_e|R_e|$, the maximum
  taken over all elements.
\end{itemize}

(I1) is not a design choice. The only pairwise non-adjacent triples among a hex's six
neighbours are the two \emph{alternating} triples, so an enemy that must be reachable
simultaneously from three mutually sealed regions is forced to present one. This is the single place where hex
geometry does real work, and it is why the construction would not transfer unchanged to a
square grid.

The local picture is worth having in mind for the whole proof (\texttt{\#} impassable,
\texttt{Z} the enemy $E_S$, \texttt{U}/\texttt{R}/\texttt{D} its three dockings---the
alternating triple; even rows drawn half a step right, the engine's convention; the three
corridors leave the box toward the three regions $R_e$, $e\in S$):

% The verbatim block below is the character-exact original from main.md and must not drift.
\begin{verbatim}
        # . #             to region R_e1
       # U #
        # Z R . .         to region R_e2
       # D #
        # . #             to region R_e3
\end{verbatim}

While $E_S$ lives, \texttt{Z} is occupied and each region's corridor is a dead end: the regions
are mutually sealed. When $E_S$ dies, \texttt{Z} becomes free and turns into a \emph{doorway} joining
the three regions. In the tight plays that decide the no-direction, however, the budget forces
every stack to be a singleton, so a dead $E_S$ was killed by three singletons standing on
\texttt{U}, \texttt{R} and \texttt{D}---and they never move again, so the doorway opens already
plugged. The confinement lemma below is the formal version of this picture.

\paragraph{Correctness, in three steps.} Under (I1)--(I4) the stack of slot $e$ can strike
exactly the $E_S$ with $S\ni e$, by one approach hex each (\textbf{reach}). Since each enemy has
3 hit points and a stack \emph{strikes} at most once per round (acting once would be the
wrong invariant under \texttt{WAIT}---Appendix~\ref{app:thm3}, Lemma~D.6), the striker sets of
distinct dead enemies are disjoint;
Lemma 3.2 then gives $3t\le\sum_e c_e\le 3q$ for $t$ kills, so $t\le q$, and at $t=q$ every
inequality is tight, forcing $3q$ singleton stacks in $3q$ slots, each striking a dead enemy,
three per enemy (\textbf{budget}). Finally---the least obvious step---a dead enemy stops
blocking its hex, so every kill opens a doorway between three regions. The \textbf{confinement}
lemma says the doorway only ever opens for stacks that have already spent their
\emph{terminal} action. The induction runs on the realized order of terminal actions; a stack
that issues \texttt{WAIT} merely postpones its terminal action to the later \texttt{WAIT}
phase, and the induction reaches it there. A dead $E_S$ was killed by exactly three singleton
stacks standing on its three dockings, which by induction are the stacks of the three slots of
$S$; so if $e\in S$, slot $e$ has already taken its terminal action. The equivalence with
\PX{} follows. The full proof, including
the general embedding lemma written out in full, is Appendix~\ref{app:thm3}.

\paragraph{The embedding.} The remaining content is that the invariants can be realized. Given
a planar incidence graph, split each element vertex of degree $d$ into a path of $d$ vertices
along the rotation order; this preserves planarity, and the cyclic order of the incident edges
becomes the linear order along the path. Then draw the resulting max-degree-3 plane graph
orthogonally on a polynomial grid. The drawing theorem needed is Tamassia--Tollis~\cite{TT89},
quoted here---as in Appendix~\ref{app:thm3}---in the form of \cite[Thm.~7.3]{DG13}: a
\textbf{connected} 4-plane graph admits an orthogonal grid drawing in $O(n^2)$ area, with at
most four bends per edge, computable in polynomial time (linear, per \cite{TT89}).
(Appendix~\ref{app:citations} records exactly what was verified from which source.) The connectivity hypothesis is not free and
splitting element vertices does not supply it: the incidence graph of $X=\{1,\dots,6\}$,
$C=\{\{1,2,3\},\{4,5,6\}\}$ is a planar yes-instance with two components. Each component is drawn
separately; the proof packs the finished \emph{boards} side by side behind impassable strips
(Appendix~\ref{app:thm3}, step 2), while the implementation packs the \emph{drawings}---two empty grid
columns between components, so $2\lambda$ hexes after scaling---which the same separation
inequality covers (the third named substitution of Appendix~\ref{app:mcscope}). Either way
the packing is
sound because the invariants (I1)--(I3) are local to a region and
(I4) only grows. Then scale by an even factor $\lambda=20$, replace a
$9\times 9$ box around each vertex by a gadget, and declare everything else impassable. The
separation inequality is explicit: gadget boxes have radius $\rho=4$, so features non-incident
in the drawing land at $L_\infty$ distance at least $\lambda-2\rho=12\ge 2$ and hex adjacency,
which requires $L_\infty$ distance 1, is impossible between them. (An earlier draft used
$\lambda=9$, giving $9-8=1$; two unrelated boxes could then touch. The repair is this
inequality.) Set-vertex boxes need a hand-built adapter, because an orthogonal drawing delivers
three edges along three of the four \emph{axis} directions while the three dockings are an
alternating triple at $120^\circ$; there are four cases, all four patterns are machine-checked,
and all four are printed in Appendix~\ref{app:adapters}.

\paragraph{What this changes.} The $O(kB)$ dynamic program of Proposition 1.1 is a statement
about a matching reach structure, not about a poor roster; Theorem 2 is not the strong endpoint
of a diversity axis but a second, independent source of hardness, with many types and a
fixed, non-adaptive reach structure---a disjoint union of 3-stars, which leaves the player no
targeting choice even though it is not a matching. Both extremes are hard. We drop the conjecture rather than
restate it.

\paragraph{The price.} Theorem 3 needs obstacles, which Theorems 1, 2 and 4 do not, and it leans
hardest on the prescribed deployment cells of Section~\ref{sec:generalized}, since the whole reduction
is carried by which slot sits where.

\subsection{Theorem 4: remove the board entirely}\label{sec:thm4}

Theorem 3 makes the reach structure as rich as planar incidence allows. The opposite extreme is
a board that imposes nothing.

\begin{thmfour}
\ALLOC{} is \textbf{strongly} NP-hard already for instances with $R=1$; \textbf{no obstacles and
no abilities of any kind}; every \emph{player} creature type of stock one, one creature per enemy stack; a
rectangular open battlefield of six rows and $4m+2$ columns; and \textbf{complete
reachability}---in the starting position, with every slot occupied, every player stack can
attack every enemy stack.
\end{thmfour}

\begin{corfourone}
The same instances are hard with the allocation \emph{given}. \PLAY{} is strongly NP-hard on
obstacle-free boards with complete reachability.
\end{corfourone}

\begin{corfourtwo}[hit-point objective]
On the same instances, replace the objective by \textbf{total enemy hit points
removed}---damage absorbed by the defence, with overkill discarded. Formally, for a play $\pi$ put
\[
  \mathrm{absorbed}(\pi):=\sum_{\text{enemy stacks }E}
    \min\bigl(\text{total damage directed at }E,\ \text{its initial pool}\bigr),
\]
and replace the question by ``is $\mathrm{absorbed}(\pi)\ge W_{\mathrm{hp}}$?'' with the
target $W_{\mathrm{hp}}=mT$. The problem remains strongly NP-hard, with the allocation free or
given.
\end{corfourtwo}

\smallskip
\noindent\emph{Proof sketch of Corollary 4.2.} Total nominal player damage is $\sum a_i=mT$ and
each enemy absorbs at most its pool $T$, so removing $mT$ hit points forces zero waste: every
stack strikes, no blow is reduced (a waiting stack that meets a \texttt{DEFEND} bonus delivers
\textbf{at most} nominal---not strictly less, since the damage formula clamps at 1 and $a_i=1$ is
legal, and the argument needs only the inequality), and every enemy absorbs exactly $T$. The
striker sets are disjoint
and $T/4<a_i<T/2$, so the tight groups are triples and form a 3-partition; conversely a
3-partition witness play delivers exactly $T$ to each enemy. Machine-checked on the same
instance families as the theorem (\src{scripts/verify\_hp\_objective.py}): the hit-point
relaxation reaches $mT$ exactly on the \THREEPART{} yes-instances, and the witness play's
absorbed total, read off the simulator, is exactly $mT$ on every one. \qedbox

\smallskip
Corollary 4.2 is owed to an external review pass, which caught an earlier draft claiming
that a hit-point objective ``would make the problem trivial''. That claim is true only where the
reach structure is a matching (Theorem 1); in general the kill-counting objective is the natural
one, but it is not the sole source of hardness.

\smallskip
\noindent\emph{Construction.} From \THREEPART{} again---NP-complete in the strong sense, as
in Theorem 2~\cite[SP15]{GJ79}: six rows, $4m+2$ columns, enemy $E_g$ with
$\mathit{att}=\mathit{def}=1$, flat damage 1, $\mathit{hp}=T$, speed 1, no flags and value 1 at
$(4g-2,3)$; player types $C_1,\dots,C_{3m}$ with $\mathit{att}=\mathit{def}=1$, flat damage
$a_i$, $\mathit{hp}=5$, no flags, value 0, stock one, and speed $s=4m+8$, deployed along the
top row at $p_j=(j-1,0)$. $R=1$, $W=m$.

\smallskip
\noindent\emph{Correctness.} The upper-bound direction uses no geometry at all: a stack takes
one terminal action per round (R9, R11), so it strikes at most once and the striker sets of
distinct enemies are disjoint, while stock one caps each blow at $a_i$ (a waiting blow
delivers at most $a_i$, which only helps this direction); killing
all $m$ enemies forces $\sum_{i\in S_g}a_i\ge T$, which sums to a 3-partition as in Theorem 2.
The geometric work is all in the yes-direction, where we choose the play. For $E_g$ at
$e_g=(X_g,3)$, $X_g=4g-2$, the three approach hexes are $q_g^1=(X_g-1,2)$, $q_g^2=(X_g,2)$,
$q_g^3=(X_g+1,3)$---genuine neighbours of $e_g$ under R1, pairwise distinct across all $3m$
seats (Lemma E.15). Row 1 is never occupied along an attack-only play (Lemma E.16), so a stack
still on $p_j=(j-1,0)$ steps into row 1, walks along it, and descends to its seat, spending at
most $w+2=4m+4\le s=4m+8$ movement points regardless of the activation order (Lemma E.17);
hence complete reachability holds with all $3m$ slots occupied, and every three-per-enemy
assignment $\varphi$ is simultaneously realizable (Lemmas E.18 and E.19). The degenerate
encodings are routed as in Theorem 1. The full proof is Appendix~\ref{app:thm124}.

\smallskip
Two honest qualifications. First, the confinement argument that keeps row 1 clear holds along
\emph{attack-only} plays---those in which every stack either passes or performs
\texttt{WALK\_AND\_\allowbreak ATTACK}. In the full model a stack may also move without attacking, wait, or
defend, and then confinement can fail. This costs nothing: confinement is used only in the
yes-direction, where the play is ours to choose, and the no-direction is geometry-free. An
earlier draft stated the lemma universally, which was a wording error. Second, the board does
impose one thing---an enemy has six neighbours, so at most six stacks can strike it in one
round (under the completed $(\star)$ no striker dies to the retaliation it draws, so no seat
is vacated mid-round and the capacity is exactly six). That bound never binds here (three per
enemy), but ``featureless'' means \emph{complete reachability plus local seat capacity}, not
``positions do not exist''.

\paragraph{Where the hardness lives, and the two-source claim.} In Theorem 4 every type has
stock one and every slot reaches every enemy, so \emph{every} injection of types into slots is
equivalent: the multiset of blows available does not depend on which slot holds which type, and
every three-per-enemy targeting is realizable from any of them. The decision that encodes the
3-partition is the choice of targets. Hence Corollary 4.1, and hence
Table~\ref{tab:where}:

\begin{table}[ht]
\centering
\begin{tabular}{@{}p{2.1cm}p{4.3cm}p{3.6cm}p{3.7cm}@{}}
\toprule
 & what forces the grouping & what the player decides & hardness \\
\midrule
Theorem 1 & board adjacency (matching reach) & how many creatures per slot & weak, and tight (Prop.~1.1) \\
Theorem 2 & board adjacency (three seats per enemy) & how to fill the seats & strong \\
Theorem 3 & planar incidence reach, single type & whom to hit & strong, allocation fixable \\
Theorem 4 & nothing & whom to hit & strong, allocation fixable \\
\bottomrule
\end{tabular}
\caption{Where the combinatorics sits in each theorem.}\label{tab:where}
\end{table}

Theorem 2 is allocation-driven: the play is forced and the allocation carries the 3-partition.
Theorems 3 and 4 are targeting-driven: the allocation is free, or fixed, and the play carries
the combinatorics. Neither collapses into the other---Theorem 2 has no targeting
choice and a rich roster, Theorem 3 a rich reach structure and a roster of one---and that is
the two-source statement of Section~\ref{sec:headline}.

\section{Machine-checking methodology}\label{sec:checking}

Every rule of Section~\ref{sec:model} was transcribed into code, every reduction was
regression-tested exhaustively on bounded instances against that transcription, and the
arithmetic was cross-checked against the shipped engine. We are deliberate about the
vocabulary: nothing here is \textbf{machine-verified} in the sense of a proof
assistant---there is no formal proof object, and the universal theorems rest on the hand
proofs; what we claim is \textbf{exhaustively regression-tested on bounded instances} and
\textbf{engine-cross-checked}, to the scope stated below; the artifact's
\src{VERIFICATION.md} carries the iteration-by-iteration record. Throughout, \emph{review
round $n$} refers to the $n$-th pass of the external review process this manuscript went
through; the numbering is internal.

\subsection{The three layers}\label{sec:layers}

\begin{enumerate}
\item \textbf{An executable transcription of the rules.} \src{scripts/homm3\_model.py}
  implements Section~\ref{sec:model} with the citations in comments. No verifier restates a
  rule; every check calls the transcription. Unit tests against hand-computed engine numbers:
  90 checks, including the \texttt{WAIT}/\texttt{DEFEND} phase mechanics the policy
  $(\ddagger)$ stands on.
\item \textbf{Exhaustive regression testing of each reduction on bounded instances.} For small
  instances we enumerate the allocation space---what that space is differs by suite:
  Theorem 1's takes every count vector summing to at most the stock, Theorem 2's every
  deployment fielding each type exactly once, the shape its witnesses take; the artifact's
  \src{README.md} itemizes it suite by suite---and, for each allocation, search the plays
  of the \textbf{attack-only fragment}, branching over targets \textbf{and} over every
  legal approach hex. That describes the Theorem 1, 2 and 3 suites; the Theorem 4 suite (and
  through it Corollary 4.2's) decides its no-direction by an arithmetic relaxation---the
  inequality of Lemma~E.20, recomputed from the instance---together with a simulation of
  every balanced assignment of stacks to enemies, and its exhaustive play search over targets
  and approach hexes runs on 2 instances $\times$ 3 allocations in the battery and on 4
  instances under \src{--full}, the smallest legal ones. The searched fragment is a
  \textbf{proper subset} of the model's play space, and the gap is closed by argument, not by search: \emph{pure movement and player
  \texttt{DEFEND}} are over-approximated---a stack that would move without attacking may
  instead \emph{vanish from the board}, which frees strictly more hexes than any move could
  and so yields a sound upper bound; \emph{\texttt{WAIT}}, which reorders terminal actions
  and leaves a blow at most nominal (Section~\ref{sec:policy}), is not searched on the
  player's side but discharged on paper, per theorem---the no-directions of Theorems 1, 2
  and 4 count blows and never refer to the acting order, and Theorem 3's confinement lemma
  inducts over the realized order of terminal actions, which covers waiting stacks at their
  postponed position; the witness plays never wait. The \emph{defence's} waiting is
  executed rather than discharged: the Theorem 1--2 suite and both Theorem 3 verifiers run
  a variant in which the enemy plays $(\ddagger)$ literally, phase by phase, with the
  \texttt{DEFEND} bonus live in the damage formula, including the published \src{def 27}
  constants. One scope note, so this is not overread: since every non-waiting blow lands
  before any postponed \texttt{DEFEND} (Section~\ref{sec:policy}) and the searched player
  never waits, no blow in these runs strikes a target that is already defending
  (instrumenting the damage calls in review round 13 recorded zero such calls---a
  measurement of record, not an invariant the battery asserts); what the $(\ddagger)$ runs
  certify is the phase machinery and that every answer is invariant under it, while the
  capped defended branch itself is exercised by the mechanics tests and by a phase-ordered
  trace in the regression battery. ``Every play'' without qualification would overstate all
  of this; which tiers are exhaustive over what, and on how many instances, is itemized in
  the artifact's \src{README.md}.
\item \textbf{Cross-checking against the shipped engine.} A C\texttt{++} harness links the
  actual VCMI battle classes---\src{DamageCalculator}, \src{CUnitState}/\src{CHealth},
  \src{CRetaliations}, and the round-queue machinery \src{battleGetTurnOrder} and
  \src{battleQueuePhase}---and prints the numbers the engine itself produces: 49 cases, 46
  exact agreements, the other 3 explained in Section~\ref{sec:floats}. Precisely what was
  cross-checked: combat arithmetic including both damage factors and the clamp at 1 (among
  them 12 damage and kill-threshold cases drawn from the reduction constructions, one of them
  a retaliation blow), the health-pool and effective-count mechanics, the retaliation-charge
  mechanics, and the single-round turn queue under
  \texttt{WAIT}/\texttt{DEFEND} states---the engine's own \src{battleGetTurnOrder}, run over
  real \src{CUnitState} units, confirms the descending \texttt{NORMAL} phase, the
  \emph{ascending} \texttt{WAIT} phase whatever the speeds (the engine half of the one-round
  lemma of Section~\ref{sec:policy}), that a defending unit leaves the queue for the round,
  and the side alternation on speed ties where our reference model documents its simplified
  tie rule. Precisely what was \emph{not}: complete battles, obstacle boards, reachability
  and approach selection, and the \texttt{DEFEND} bonus's server-side arithmetic and expiry
  (\src{BattleActionProcessor::\allowbreak doDefendAction}, \src{BattleInfo::\allowbreak nextTurn}),
  which cannot be driven without the full game server and are instead quoted from source
  with file-line citations and regression-tested in the reference model. No constructed
  instance has been played inside VCMI.
\end{enumerate}

\subsection{One caveat about arithmetic}\label{sec:floats}

The damage formula carries a defence-factor cap at $0.7$ under a floor, and our reductions
are positioned so that nothing they consume depends on how that arithmetic is evaluated.
Theorems 1, 2 and 4 set $\Delta=0$, so their witness plays form no fractional product at
all; their no-directions consume only the bound ``delivered at most nominal''
(Lemmas E.5, E.12 and E.20, for Theorems 1, 2 and 4 respectively), true under every
arithmetic (for Theorem 1 this is a statement about the
model's exact integer arithmetic: its instances carry hit points in binary, beyond what the
engine's \src{int32} stack counts could represent, so no engine-agreement claim is made in
that regime); Theorem 3's undefended blows sit at $0.025\cdot 26=0.65$, clear of the cap,
and its resource lemma is stated for every $\mu\in(0,1)$, so it holds on either side of the
cap that the defended branch (Section~\ref{sec:policy}) does cross. Two independent one-ULP
divergences from the exact rationals are known and documented, and no theorem of this paper
depends on either: VCMI's hand-written JSON parser loads the cap literal $0.7$ as
$0.7000000000000001$, one ULP above the correctly rounded double, so the engine's
\src{std::floor} discards a point of damage whenever $\text{base}\times 0.3$ is an exact
integer---3 of the 49 engine cases, all with defence far above attack; and the reference
simulator evaluates the formula in Python floats, which part from the exact rationals in the
opposite direction on a different constant ($1-0.025\cdot 12$ is the double nearest $7/10$,
below it). On the constants of the constructions the two arithmetics agree on every stack
size the proofs reason about: the battery checks equality up to $c=2000$ creatures on the
$(\star)$ waiting blow, on both Theorem 3 branches at the historical constants and on the
defended branch at the published ones, and pins the first split of the published undefended
branch at $c=180$---beyond every stack the Theorem 3 search forms, $3q\le 12$
(Appendix~\ref{app:thm3}). On the empirical certification suites the two split on 202 of
909638 damage calls, each a float result one point below the exact one, none on a branch
that decides a certified number. The details, the engine's own reading of its constant, and
the cross-check that re-derives every certificate of the empirical companion note under
exact arithmetic are in the artifact (\src{engine-check/REPORT.md},
\src{empirics/scripts/exact\_\allowbreak arithmetic\_\allowbreak crosscheck.py}).

\subsection{What the checks caught}\label{sec:errors}

The checks caught three substantive errors in earlier versions of this work---two in proofs
believed correct and written out in full (the first version of Theorem 1 used an
\emph{attacking} garrison, whose retaliation strikes the same enemy a second time in the
same round, so the reduction decided \SUBSETSUM{} with budget $2B$; the first version of
Theorem 2 used a one-row battlefield, where a third stack cannot reach its enemy past an
ally) and one in the checking apparatus itself (the featureless verifier ran the attacking
defence under both policy labels, so the player stacks retaliated and three kills were
reported where the arithmetic admits one). External review caught errors that bounded
checks cannot see, among them a false proposition (an earlier statement of Proposition 1.1
lacked its one-creature-per-enemy-stack hypothesis; the DP suite now plays the refuting
instance as a negative control), a claim about instances nobody had constructed (false by
Corollary 4.2), the side-column simplification of Section~\ref{sec:generalized}, and a
sentence in prose about what a table showed (the empirical companion note carries the
retraction). Two further errors---in the empirical study's exact solver and in its
prompting protocol---belong to the companion note and are recorded with its artifact
(\src{empirics/RESULTS.md}); an earlier version of this section counted them among the
errors the checks caught, and one of them was itself found by external review. The full
accounts, error by error and with who found each, are the artifact's
\src{VERIFICATION.md}.

\subsection{Reproducing}\label{sec:reproducing}

% alltt rather than verbatim: the block is set small enough that no line wraps at the margin.
\begingroup\scriptsize
\begin{alltt}
python3 scripts/test_regressions.py          # battery: every counter pinned, cheap suites re-run
python3 scripts/verify_x3c.py --full --vacate  # Theorem 3, q \(\le\) 4 corpus (minutes; not re-run above)
python3 scripts/crosscheck_sol.py --full     # Theorem 3 at the PUBLISHED def 27, hp 4, \(\mu\) = 0.35
cd engine-check && ./build.sh && ./run.sh && python3 compare.py   # needs a VCMI checkout
\end{alltt}
\endgroup

There are no dependencies beyond the Python standard library, except for the engine
cross-check. Every suite the battery runs or pins---the command that runs it, what it runs
at what scale, and its current outcome---is a row of the table in the artifact's
\src{README.md}, rendered from \src{verification\_\allowbreak manifest.json} rather than
maintained by hand; the tiers the battery neither runs nor pins (the unrestricted
\src{--full} tier of the Theorem 1--2 suite, the pipeline fuzz, the skip classifier) are
listed with their commands above that table. \src{test\_\allowbreak regressions.py}
re-renders the table in check mode and pins every counter it quotes to the artifact that
generates it, in five ways, itemized suite by suite in the \src{README.md}: by re-running the
generating suite and parsing its own final line (the mechanics, geometry, reduction and
embedding suites, seconds to minutes each); by re-deriving the engine verdicts offline from
the shipped engine outputs; by the lengths of the shipped empirical files for the instance
and response counts; by a recorded log plus an in-process replay of its witness half for the
arithmetic cross-check---the split count of Section~\ref{sec:floats} is a pin to that log,
not a re-run, since the cross-check re-runs three scoring suites twice and stays outside the
battery by the same policy as every other multi-minute tier; and by sweeping the documents
that quote the two \src{--full} Theorem 3 tiers, which it does not re-run. It also sweeps
the counters this paper and the companion note quote, and bans retired claims verbatim;
what that guard does and does not establish is recorded in the artifact's
\src{VERIFICATION.md}.

\paragraph{Data and code availability.}
Everything this section runs is published as an artifact repository~\cite{Par26}, release
tag \src{v1.5}---cite and check out the tag, not the moving branch;
\src{scripts/check\_\allowbreak artifact\_\allowbreak repo.py} pins the tag to the paper: every
path the paper or the companion note names (their count is pinned in the manifest), the
manifest counters, and byte-identical copies of \src{main.md} and of the companion note
(\src{paper/companion-\allowbreak empirics.md}). All file paths in this paper are relative to the
artifact root. The engine cross-check was run against a VCMI~\cite{VCMI26} checkout at
commit \src{b5cee70}; every VCMI file this paper cites is byte-for-byte identical to the
public tree at commit \src{deeab240} (\src{develop}, 2026-06-19), which is the anchor a
reader should check out.

\section{Related work}\label{sec:related}

\paragraph{Turn-based tactics.} Gao~\cite{Gao19} gives the only prior complexity study of a
commercial turn-based tactical RPG, proving a simplified Fire Emblem PSPACE-complete and its
round-bounded version NP-complete. Two features matter for the comparison: every numeric
attribute in that model is a constant bounded by 8---so the hardness is purely combinatorial,
and the NP-completeness trivially strong---and the player's army is \emph{given}, with no
allocation decision anywhere. That hardness comes from geometry; ours (in Theorems 1, 2 and 4)
comes from arithmetic.

\paragraph{Attrition games.} The closest prior result is Furtak and Buro~\cite{FB10}, who study
attrition games on graphs---nodes are units with $\langle\text{health},\text{attack}\rangle$,
edges say who may attack whom, movement removed. Their Theorem 7 shows that \emph{attack
partitioning}, in which a unit divides its attack power among its targets, is NP-hard by
reduction from \SUBSETSUM; their Theorem 8 partitions a set of attacking units against health
thresholds. \textbf{These are the mechanisms of our Theorems 1 and 2, published in 2010, and we
do not claim them as new.} Our claims against \cite{FB10} are narrower: the resource we split is
a native player-facing action of a shipped game rather than a modelling device introduced to
interpolate between models; it is a discrete count of creatures rather than a divisible scalar;
the map from allocation to reachable enemy is realized by board geometry rather than supplied as
an abstract graph---though, as Section~\ref{sec:generalized} concedes, we choose the deployment
cells, so this map is not \emph{forced} by the shipped formation tables either; and our hardness
is strong, whereas every numeric hardness result in \cite{FB10} is weak. A systematic sweep of AIIDE (2005--2025, workshops included) and CIG/CoG
(2008--2025)---roughly 2900 titles over two independent passes---found no other hardness result
about combat or force allocation at either venue. \cite{FB10} closes with the question our
Theorem 1 sits inside: what is the smallest $k$ for which the $k$ vs.\ $n$ problem is NP-hard?

\paragraph{Card games.} Hoffmann, Lynch and Winslow~\cite{HLW20} reduce \THREEPART{} to
board-scaled \emph{Lethal} in Hearthstone: $3n$ minions with attack $4a_i$ against $n$ taunts of
health $4S/n$, no overkill possible, winning assignments exactly the 3-partitions. \textbf{This
is the closest published relative of our Theorems 2 and 4, in a shipped commercial game, six
years earlier.} Their Theorem 3.1 states NP-hardness; since the reduction is from \THREEPART{}
on a board-scaled instance it is strong, and strong NP-hardness as such is not new in this
genre---we do not claim it is. Bosboom and Hoffmann~\cite{BH17} prove Netrunner mate-in-1 weakly NP-hard from
2-\textsc{Partition}; Rom\~ao et al.~\cite{RPU25} model an aggressive Flesh and Blood turn as an
ILP containing 0-1 Knapsack, obtaining a weak hardness result that the authors themselves
expect to be pseudo-polynomially solvable. UNO is NP-hard for a single player~\cite{DDHUUU14};
perfect-information Hearthstone is PSPACE-hard~\cite{Zha23}; Magic: The Gathering is
Turing-complete~\cite{CBH19}.

\paragraph{Classical ancestry.} With a single type, $R=1$, and each slot facing one enemy, our
problem is 0-1 Knapsack~\cite{Kar72} and, with $v_j=a_j$, \PARTITION~\cite[SP12]{GJ79}. With
unit-multiplicity heterogeneous attackers and per-slot thresholds it is a form of \textsc{Bin
Covering}, strongly NP-hard by Assmann, Johnson, Kleitman and Leung~\cite{AJKL84}.
Weapon-Target Assignment is NP-complete~\cite{LW86} but draws its hardness from a nonlinear
probabilistic objective, so it is an ancestor rather than a competitor. The Colonel Blotto
literature runs the other way: closed-form equilibria for the discrete game~\cite{Har08} and
polynomial-time equilibrium computation for the general case~\cite{ADHLMS19}; hardness appears
only when the resource stops being homogeneous~\cite{DSST21}.
Stripped of the game, Theorem 4 \emph{is} bin covering, and we say so: the interest is in how
little of HoMM3 the hardness needs, not in the combinatorial core.

\paragraph{Composition rather than allocation.} Ponomarenko and Sirotkin~\cite{PS20} prove
optimal team choice in the auto-battler Dota Underlords NP-complete. That is the exact complement
of our framing: there the difficulty is entirely in \emph{which} units to field and the
arrangement is free; in our Theorem 1 there is nothing to choose and the difficulty is entirely
in the split.

\paragraph{HoMM3.} There is no prior complexity result for any title in the Heroes of Might and
Magic, King's Bounty, Disciples, Age of Wonders or Master of Magic series. We are aware of two
academic papers on HoMM3, neither about complexity: Diochnos~\cite{Dio10} models
the random secondary-skill offers of levelling up; Kowalski et al.~\cite{KMPPPS18} generate
balanced maps from terrain features.

\section{Open problems}\label{sec:open}

We list these in the order we would attack them, and we deliberately defer several that a reader
might expect.

\begin{enumerate}
\item \textbf{A natural victory objective.} All four theorems use the artificial deadline $R=1$
  and the objective ``value destroyed''. Replacing it by ``eventually eliminate the defence'' is
  the most valuable single improvement. A promising route makes every item creature a shooter
  with one shot, separated from the defence by an impassable barrier, so that surviving enemies
  can never be damaged after the first volley; it needs shooting and ammunition in the executable
  model, which the formal model already retains.
\item \textbf{Fixed $k$.} The shipped game has $k=7$; all four constructions need $k$ to grow. We
  conjecture the problem is in P for fixed $k$ and $R=1$. Proving it would give the contrast
  result the separation wants.
\item \textbf{Native deployment formations.} All four theorems take the slot-to-hex map as input
  (Section~\ref{sec:generalized}). Recovering any of them when deployment comes from the engine's
  formation tables is open, and is the sharpest form of the objection ``your board is not a game
  position''.
\item \textbf{Single type on a featureless board.} Theorem 3 needs obstacles; Theorem 4 needs a
  diverse roster. Whether complete reachability plus a single type is tractable is open even
  with one creature per enemy stack: the obvious cardinality-constrained knapsack is only an
  upper bound, because its achievability needs the chosen stacks to hold seats at distinct
  enemies \emph{simultaneously}, and complete reachability---a property of the starting
  position---does not supply that (a six-hex instance in the executable model has knapsack
  value 2 and true optimum 1; Theorem 4 closes the same gap with the routing argument of
  Appendix~\ref{app:thm124}). Enemy stacks of many creatures, where partial kills score, may
  be harder still, and that is the version a referee is most likely to ask for.
\item \textbf{Bounded speed.} Theorem 3 sets creature speed to the board size to avoid equalizing
  corridor lengths. Legal, but inelegant.
\item \textbf{Adversarial defence.} The NP-membership argument given here does not extend to an
  optimizing opponent, and we supply no lower bound for that variant: its complexity is open.
  The metatheorems of de Haan and Wolf~\cite{dHW18} suggest the second
  level of the polynomial hierarchy rather than PSPACE if one player is strategically restricted.
\item \textbf{Approximation.} None of our results gives inapproximability. Since Theorem 4 is bin
  covering, positive approximation results would have to beat what is already known offline, so
  this is less attractive than it looks.
\end{enumerate}

\section{Conclusion}\label{sec:conclusion}

Planning one round of HoMM3 combat is hard in two separable ways: choosing how to divide an army
among slots, and choosing what to hit. Each remains strongly NP-hard under conditions that remove
the other---a fixed single-type army for the second, a board that dictates the targeting for the
first---and the tractable case we exhibit is tractable because its reach structure is a perfect
matching over single-creature stacks, not because its roster is poor. Every rule used is cited to a line of an engine that
reproduces the shipped game's combat, and every reduction has been regression-tested on bounded
instances against an executable transcription of those rules, to the scope itemized in
Section~\ref{sec:checking}---a process that changed two of the proofs in this paper.

\section*{Acknowledgements}\label{sec:ack}

The proofs, the machine-checking apparatus and the drafts of this paper were produced with
substantial AI assistance: several large language models worked as separate agents on the model
transcription, proof exploration and drafting, the verification code, the literature survey and
the empirical harness. Concretely: the model transcription, the proofs, the verification code
and the drafts were written by Anthropic's Claude agents (Claude Fable 5 and 5.1, Claude Opus 5 and 5.5,
run through Claude Code); the second, independent proof of Theorem 3 by an OpenAI GPT-5.6 agent
(\src{codex}); and the adversarial review rounds by OpenAI GPT-5.6 and GPT-6 (\src{codex} CLI,
one round through ChatGPT Pro), DeepSeek v4-pro, and Claude Fable 5, Fable 5.1, Opus 5 and
Opus 5.5 agents. Theorem 3 was proved twice, independently and in parallel, by two agents
that were not permitted to see each other's work; the two proofs agreed and were merged. The
author treats that agreement as an error-catching redundancy, not as independent scientific
validation. The mechanics and the constructions were tested to the scope stated precisely in
Section~\ref{sec:checking}---bounded instances and selected engine mechanics, not all statements
and not a formal proof object. Every reference cited was independently located, with the one exception recorded in
Appendix~A; the depth to
which each load-bearing statement from the literature was verified varies by reference and is
recorded in Appendix~A---the two dependencies of Theorem 3 in detail, and every citation not
read in full listed explicitly. The subjects of the empirical companion note
(\src{paper/companion-empirics.md}) are models of the same family that assisted in writing
this paper; the note carries that conflict statement, and this paper claims nothing about
its outcome. The author reviewed the text and is responsible for the content.

\begin{thebibliography}{DDHUUU14}

\bibitem[ADHLMS19]{ADHLMS19} A.~Ahmadinejad, S.~Dehghani, M.~Hajiaghayi, B.~Lucier, H.~Mahini,
S.~Seddighin. From Duels to Battlefields: Computing Equilibria of Blotto and Other Games.
\emph{Mathematics of Operations Research} 44(4):1304--1325, 2019.
\url{https://doi.org/10.1287/moor.2018.0971}

\bibitem[AJKL84]{AJKL84} S.~F. Assmann, D.~S. Johnson, D.~J. Kleitman, J.~Y.-T. Leung.
On a dual version of the one-dimensional bin packing problem.
\emph{Journal of Algorithms} 5(4):502--525, 1984.
\url{https://doi.org/10.1016/0196-6774(84)90004-X}

\bibitem[BH17]{BH17} J.~Bosboom, M.~Hoffmann. Netrunner Mate-in-1 or -2 is Weakly NP-Hard.
arXiv:1710.05121, 2017. \url{https://arxiv.org/abs/1710.05121}

\bibitem[BM04]{BM04} J.~M. Boyer, W.~J. Myrvold. On the Cutting Edge: Simplified $O(n)$ Planarity
by Edge Addition. \emph{Journal of Graph Algorithms and Applications} 8(3):241--273, 2004.
\url{https://doi.org/10.7155/jgaa.00091}

\bibitem[CBH19]{CBH19} A.~Churchill, S.~Biderman, A.~Herrick. Magic: The Gathering is Turing
Complete. arXiv:1904.09828, 2019; FUN 2021, LIPIcs 157, art.~9, pp.~9:1--9:19.
\url{https://doi.org/10.4230/LIPIcs.FUN.2021.9}

\bibitem[DDHUUU14]{DDHUUU14} E.~D. Demaine, M.~L. Demaine, N.~J.~A. Harvey, R.~Uehara, T.~Uno,
Y.~Uno. UNO is hard, even for a single player. \emph{Theoretical Computer Science}
521:51--61, 2014. \url{https://doi.org/10.1016/j.tcs.2013.11.023}

\bibitem[DF86]{DF86} M.~E. Dyer, A.~M. Frieze. Planar 3DM is NP-complete.
\emph{Journal of Algorithms} 7(2):174--184, 1986.
\url{https://doi.org/10.1016/0196-6774(86)90002-7}

\bibitem[DG13]{DG13} C.~A. Duncan, M.~T. Goodrich. Planar Orthogonal and Polyline Drawing
Algorithms. In R.~Tamassia (ed.), \emph{Handbook of Graph Drawing and Visualization}, ch.~7,
pp.~238--261, CRC Press, 2013. \url{https://doi.org/10.1201/b15385-10}

\bibitem[dHW18]{dHW18} R.~de Haan, P.~Wolf. Restricted Power --- Computational Complexity Results
for Strategic Defense Games. FUN 2018, LIPIcs 100, art.~17, pp.~17:1--17:14.
\url{https://doi.org/10.4230/LIPIcs.FUN.2018.17}

\bibitem[Dio10]{Dio10} D.~I. Diochnos. Leveling-Up in Heroes of Might and Magic III.
FUN 2010, LNCS 6099, pp.~145--155. \url{https://doi.org/10.1007/978-3-642-13122-6_16}

\bibitem[DMP64]{DMP64} G.~Demoucron, Y.~Malgrange, R.~Pertuiset. Graphes planaires: reconnaissance
et construction de repr\'esentations planaires topologiques. \emph{Revue Fran\c{c}aise de
Recherche Op\'erationnelle} 8:33--47, 1964.

\bibitem[DSST21]{DSST21} S.~Dehghani, H.~Saleh, S.~Seddighin, S.-H. Teng. Computational Analyses
of the Electoral College: Campaigning Is Hard But Approximately Manageable. AAAI 2021, pp.~5294--5302.
\url{https://doi.org/10.1609/aaai.v35i6.16668}

\bibitem[FB10]{FB10} T.~Furtak, M.~Buro. On the Complexity of Two-Player Attrition Games Played on
Graphs. AIIDE 2010, pp.~113--119. \url{https://doi.org/10.1609/aiide.v6i1.12410}

\bibitem[Gao19]{Gao19} J.~Gao. The Computational Complexity of Fire Emblem Series and similar
Tactical Role-Playing Games. arXiv:1909.07816, 2019. \url{https://arxiv.org/abs/1909.07816}

\bibitem[GJ79]{GJ79} M.~R. Garey, D.~S. Johnson. \emph{Computers and Intractability}.
Freeman, 1979.

\bibitem[Har08]{Har08} S.~Hart. Discrete Colonel Blotto and General Lotto games.
\emph{International Journal of Game Theory} 36(3--4):441--460, 2008.
\url{https://doi.org/10.1007/s00182-007-0099-9}

\bibitem[HLW20]{HLW20} M.~Hoffmann, J.~Lynch, A.~Winslow. Mad Science is Provably Hard: Puzzles in
Hearthstone's Boomsday Lab are NP-hard. arXiv:2010.08862, 2020.
\url{https://arxiv.org/abs/2010.08862}

\bibitem[HT74]{HT74} J.~Hopcroft, R.~Tarjan. Efficient Planarity Testing. \emph{Journal of the
ACM} 21(4):549--568, 1974. \url{https://doi.org/10.1145/321850.321852}

\bibitem[Kar72]{Kar72} R.~M. Karp. Reducibility Among Combinatorial Problems.
In \emph{Complexity of Computer Computations}, Plenum Press, 1972, pp.~85--103.
\url{https://doi.org/10.1007/978-1-4684-2001-2_9}

\bibitem[KMPPPS18]{KMPPPS18} J.~Kowalski, R.~Miernik, P.~Pytlik, M.~Pawlikowski, K.~Piecuch,
J.~S\k{e}kowski. Strategic Features and Terrain Generation for Balanced Heroes of Might and
Magic III Maps. IEEE CIG 2018, pp.~1--8.
\url{https://doi.org/10.1109/CIG.2018.8490430}

\bibitem[LW86]{LW86} S.~P. Lloyd, H.~S. Witsenhausen. Weapons allocation is NP-complete.
\emph{Proc. 1986 Summer Computer Simulation Conference}, Reno, NV, pp.~1054--1058, 1986.

\bibitem[Par26]{Par26} I.~Parfenchuk. Artifact repository for this paper: model
transcription, proofs, verification scripts, engine cross-check and empirical harness.
\url{https://github.com/uson1x/homm3-hardness}, release tag \src{v1.5}, 2026.

\bibitem[PS20]{PS20} A.~A. Ponomarenko, D.~V. Sirotkin. Dota Underlords game is NP-complete.
arXiv:2007.05020, 2020. \url{https://arxiv.org/abs/2007.05020}

\bibitem[RPU25]{RPU25} L.~G. Rom\~ao, S.~P. de Paula, E.~T. Ueda. Optimizing for
aggressive-style strategies in Flesh and Blood is NP-hard. arXiv:2501.11683, 2025.
\url{https://arxiv.org/abs/2501.11683}

\bibitem[TT89]{TT89} R.~Tamassia, I.~G. Tollis. Planar grid embedding in linear time.
\emph{IEEE Transactions on Circuits and Systems} 36(9):1230--1234, 1989.
\url{https://doi.org/10.1109/31.34669}

\bibitem[VCMI26]{VCMI26} The VCMI Project. VCMI: open-source engine reimplementation of
Heroes of Might and Magic III. \url{https://github.com/vcmi/vcmi}, GPLv2. Cited files
anchored at public commit \src{deeab240} (\src{develop} branch, 2026-06-19).

\bibitem[Zha23]{Zha23} Z.~Zhang. Perfect Information Hearthstone is PSPACE-hard.
arXiv:2305.12731, 2023. \url{https://arxiv.org/abs/2305.12731}

\end{thebibliography}

\appendix

\section{Status of the citations}\label{app:citations}

The two literature dependencies of Theorem 3 were checked on 2026-08-03; here is exactly what
was verified and from what.

\begin{itemize}
\item \textbf{\cite{DF86}}---\textbf{read in full against the published paper} (scan of
  \emph{J. Algorithms} 7:174--184). Planarity is defined on the triple/element incidence graph:
  ``We have a vertex for each element \dots\ and each triple \dots\ There is an edge connecting a
  triple to an element if and only if the element is a member of the triple. \dots\ We will say
  that the instance is \emph{planar} if G is planar'' (p.~175). \textbf{Lemma 2.2 states ``Planar
  X3C is NP-complete'' outright}, proved directly from Planar 1-3SAT, so the earlier ``Planar
  3DM, hence Planar X3C by identity'' detour is unnecessary and has been removed. Bonus, from
  p.~178: their X3C instances have every element in two or three sets, so \PX{} remains
  NP-complete under that degree restriction.
\item \textbf{\cite{TT89}}---the exact statement is quoted from an authoritative secondary
  source, the \emph{Handbook of Graph Drawing and Visualization} chapter on orthogonal drawings
  \cite{DG13} (hosted by Tamassia), whose Theorem 7.3 credits \cite{TT89} with: a biconnected
  4-plane graph admits an orthogonal grid drawing in $O(n^2)$ area with at most $2n+4$ bends,
  and \textbf{a connected 4-plane graph, with at most $2.4n+2$ bends and no edge bending more
  than four times}. Theorem 7.3 itself states no running time; the linear-time claim is the
  title of \cite{TT89} and the chapter's framing text, and is corroborated by the primary
  abstract quoted below. The argument uses only (D5), polynomial time. The connected case is
  the one Lemma D.4 uses (its graph is connected after step 2, with maximum degree 3). The abstract of the IEEE original was additionally
  obtained as recorded in a bibliographic index (the OpenAlex record, reconstructed from the
  record's inverted word index---not the publisher's own text) and states verbatim: an
  $O(n)$-time algorithm producing grid embeddings with ``the total number of bends is at most
  2.4n+2'', ``the number of bends along each edge is at most 4'', edge length $O(n)$, and area
  $O(n^2)$---every number the lemma relies on thus matches an independent record of the
  primary abstract. The abstract does not state the connectivity hypothesis, so the
  \textbf{connected-case theorem itself---the statement Lemma D.4 uses---still rests on the
  Handbook alone} (the constant $2.4n+2$ is incidental: the proof never uses it); the paywalled
  proof body and the biconnected $2n+4$ refinement (not used here) remain unread.
\end{itemize}

Everything else cited was read in full, with the following exceptions, listed so that the
previous sentence cannot silently overclaim. \cite{LW86} could not be located: the 1986
proceedings are not online, and the entry follows the standard secondary citation; it is cited
only for the attribution of NP-completeness. \cite{GJ79} was consulted for its catalogue
entries SP12 and SP15, not read cover to cover. \cite{ADHLMS19}, \cite{HT74}, \cite{BM04} and
\cite{DMP64} are cited for attribution of known algorithms; their abstracts and theorem
statements were consulted, not the full texts. \cite{HLW20} was read in detail in its sections 1--3 and
its load-bearing theorem statements were transcribed verbatim, but not cover to cover;
\cite{PS20} was verified from its abstract and problem statement; \cite{BH17} from its
abstract and introduction. These last three are cited as related work---for which game they study
and that a hardness result exists---and none of their constructions is reused anywhere in
this paper.

\section{Rule-to-source table}\label{app:rules}

Table~\ref{tab:rules} gives every rule tag used in the body, with its \src{file:line} in the
VCMI checkout at commit \src{b5cee70}. Paths are relative to the repository root;
\src{lib/battle/} is abbreviated \src{lb/} and \src{server/battles/} is abbreviated
\src{sb/}.

\begingroup
\small
\begin{longtable}{@{}p{0.8cm}p{6.0cm}p{6.8cm}@{}}
\caption{Rule tags, their statements, and their sources in the VCMI checkout.}\label{tab:rules}\\
\toprule
tag & rule & source \\
\midrule
\endfirsthead
\toprule
tag & rule & source \\
\midrule
\endhead
\bottomrule
\endfoot
R1 & hex grid, offset coordinates, six neighbour directions & \src{lb/BattleHex.h:60-73, 147-178} \\
R2 & hex distance via axial coordinates & \src{lb/BattleHex.h:195-210} \\
R3 & impassable hexes: battlefield set, obstacle objects, siege walls & \src{lb/CBattleInfoCallback.cpp:1328-1391} \\
R4 & a living unit blocks its hex; a dead one does not & \src{lb/CBattleInfoCallback.cpp:1355-1360} \\
R5 & stack health pool \src{(fullUnits, firstHPleft)} & \src{lb/CUnitState.cpp:183-186} \\
R6 & effective count \src{fullUnits + [firstHPleft > 0]} & \src{lb/CUnitState.cpp:282-285} \\
R7 & damage: attack/defence factors and the lower clamp at 1 & \src{lb/DamageCalculator.cpp:210-224, 322-337, 123-131, 576-577} \\
R8 & kill threshold; overkill discarded & \src{lb/DamageCalculator.cpp:522-531}; \src{lb/CUnitState.cpp:202-203} \\
R9 & turn order: speed desc., side, slot; \texttt{WAIT} phase after \texttt{NORMAL}, ascending speed & \src{lb/BattleInfo.cpp:978-1006}; \src{lb/CBattleInfoCallback.cpp:496-509, 601-623}; phase enum \src{lb/Unit.h:33-43} \\
R10 & speed doubles as initiative and movement range & \src{lb/CUnitState.cpp:589-600} \\
R11 & movement: unweighted BFS over enterable hexes; \texttt{WALK\_AND\_ATTACK} & \src{lb/CBattleInfoCallback.cpp:1411-1469}; \src{sb/BattleActionProcessor.cpp:216-352} \\
R12 & retaliation after the blow; once per round; the dead do not retaliate & \src{sb/BattleActionProcessor.cpp:298-334}; \src{lb/CUnitState.cpp:484-490} \\
R13 & \texttt{DEFEND}: ends turn, $+20\,\%$ defence (floor $+1$) until next turn & \src{sb/BattleActionProcessor.cpp:160-212, 693}; turn passes on: \src{sb/BattleFlowProcessor.cpp:804-868} \\
R14 & shipped bounds: $11\times 17$ board, $\le 7$ slots, formation tables for deployment & \src{lb/BattleHex.h:19-24}; \src{lib/constants/NumericConstants.h:32}; \src{config/gameConfig.json:625, 635-643, 653-661} \\
R15 & JSON parser loads the literal $0.7$ as $0.7000000000000001$ & \src{lib/json/JsonParser.cpp:536-551} \\
R16 & first and last columns are not usable cells (\src{SIDE\_COLUMN}) & \src{lb/BattleHex.h:97-100}; \src{lb/CBattleInfoCallback.cpp:1321-1326} \\
R17 & equal-initiative ties alternate sides by \src{sideThatLastMoved}; attacker priority only on the first turn & \src{lb/CBattleInfoCallback.cpp:474-509} \\
\end{longtable}
\endgroup

\section{The four adapter patterns of Lemma D.3}\label{app:adapters}

An orthogonal drawing delivers the three edges of a set-vertex along three of the four axis
directions; the three dockings are the alternating triple
$\{\texttt{TOP\_LEFT},\texttt{RIGHT},\texttt{BOTTOM\_LEFT}\}$. The gap is bridged by a
$9\times 9$ box, one pattern per missing axis direction. In the local frame the enemy \texttt{Z}
sits at $(4,4)$ with row 4 even; ports (where a corridor crosses the box boundary) are the four
boundary midpoints. Notation: \texttt{\#} impassable, \texttt{.} free corridor hex, \texttt{Z}
the enemy, and $\texttt{U}=(4,3)$, $\texttt{R}=(5,4)$, $\texttt{D}=(4,5)$ the three dockings;
even rows are drawn half a step to the right, as on the hex board---from an even row the
two upper neighbours sit at columns $x$ and $x+1$ (rule R1). Each arm is a path from its
port to its docking; no two arms share or touch a hex; the arm touches \texttt{Z} only at the
docking; the unused port stays impassable. All four patterns are machine-checked hex by hex
(\src{verify\_x3c.py::check\_enemy\_adapters}).

% The verbatim blocks below are character-exact copies of main.md and must not drift.

\noindent Missing \texttt{LEFT} (all three arms straight):

\begin{verbatim}
 # # # # . # # # #
# # # # . # # # #
 # # # # . # # # #
# # # # U # # # #
 # # # # Z R . . .
# # # # D # # # #
 # # # # . # # # #
# # # # . # # # #
 # # # # . # # # #
\end{verbatim}

\noindent Missing \texttt{RIGHT} (the left arm dips one row to reach \texttt{D}; the bottom arm
swings around the box to reach \texttt{R}, preserving the cyclic order):

\begin{verbatim}
 # # # # . # # # #
# # # # . # # # #
 # # # # . # # # #
# # # # U # # # #
 . . . # Z R . # #
# # . . D # . # #
 # # # # # # . # #
# # # # . . . # #
 # # # # . # # # #
\end{verbatim}

\noindent Missing \texttt{BOTTOM} (left arm dips to \texttt{D}; top and right arms straight):

\begin{verbatim}
 # # # # . # # # #
# # # # . # # # #
 # # # # . # # # #
# # # # U # # # #
 . . . # Z R . . .
# # . . D # # # #
 # # # # # # # # #
# # # # # # # # #
 # # # # # # # # #
\end{verbatim}

\noindent Missing \texttt{TOP} (left arm rises to \texttt{U}; right and bottom arms straight):

\begin{verbatim}
 # # # # # # # # #
# # # # # # # # #
 # # # # # # # # #
# # . . U # # # #
 . . . # Z R . . .
# # # # D # # # #
 # # # # . # # # #
# # # # . # # # #
 # # # # . # # # #
\end{verbatim}

\section{The full proof of Theorem 3}\label{app:thm3}

This appendix makes Theorem 3 and Corollary 3.1 self-contained: the complete construction, the
embedding lemma with its proof, and the correctness argument, expanding the outline of
Section~\ref{sec:thm3}. The statements of Theorem 3, Corollary 3.1 and the resource lemma
(Lemma 3.2), the four invariants (I1)--(I4), and the enemy-gadget picture are in
Section~\ref{sec:thm3} and are not repeated; the adapter patterns are printed in
Appendix~\ref{app:adapters}. Citation provenance for the two literature
dependencies---Dyer--Frieze~\cite{DF86} and Tamassia--Tollis~\cite{TT89} via
\cite{DG13}---is Appendix~\ref{app:citations}.

\subsection{The instance}

Given a \PX{} instance $(X,C)$ with $|X|=3q$---a universe $X$, a collection $C$ of 3-element
subsets whose bipartite incidence graph $G(X,C)$ is planar, asking for $q$ pairwise disjoint
members covering $X$---build the \ALLOC{} instance $G_3(X,C)$:

\begin{itemize}
\item \textbf{Player army.} One creature type $P$ with $\mathit{att}=1$, $\mathit{def}=1$, flat
  damage $1$, $\mathit{hp}=4$, $\mathit{spd}=\sigma$ (the board's hex count),
  $\mathit{value}=0$, and stock exactly $3q$. There are $k=3q$ slots, one per element $e\in X$,
  at deployment hexes $p_e$.
\item \textbf{Defence.} For each $S\in C$, one stack $E_S$ of \textbf{one} creature of type $Q$
  with $\mathit{att}=1$, $\mathit{def}=27$, flat damage $1$, $\mathit{hp}=3$, $\mathit{spd}=1$,
  $\mathit{value}=1$, at hex $z_S$, playing the policy $(\ddagger)$ of Section~\ref{sec:policy}.
\item \textbf{Board.} Produced by Lemma D.4 below, satisfying (I1)--(I4). \textbf{Question:}
  $R=1$, $W=q$.
\end{itemize}

With $\Delta=\mathit{att}(P)-\mathit{def}(Q)=-26$ and $0.025\cdot 26=0.65<0.7$, an
\emph{undefended} blow never reads the cap constant, and a stack of $c$ player creatures
delivers nominal damage $D(c)=\max(1,\lfloor 0.35c\rfloor)$ (Section~\ref{sec:thm3}); the
defended branch does read it---see below. Every numeric parameter is a
constant except $\sigma$, the stock $3q$ and $W=q$, all bounded by a polynomial in $|X|+|C|$, so
the instance is polynomial even under \textbf{unary} encoding---which is what makes the hardness
strong. The choice $\mathit{def}=27$ is not delicate: any $\mathit{def}(Q)>\mathit{att}(P)$ gives
$\mu\in(0,1)$, which is all Lemma 3.2 needs; $27$ merely keeps $0.65$ clear of the
engine's cap hazard of Section~\ref{sec:floats} on both sides. (The reference simulator's own
float hazard, also in Section~\ref{sec:floats}, is not dodged by the constant: $1-0.65$ is the
double nearest $7/20$, which lies below it, so $\lfloor c\cdot 0.35\rfloor$ first loses a point
at $c=180$, i.e.\ $q\ge 60$---far beyond every searched instance, $3q\le 12$, and immaterial
to the lemma, which only needs $\mu\in(0,1)$; the battery pins both facts.) A \emph{defending} $E_S$ has
defence $27+\lfloor 27\cdot 20/100\rfloor=32$, so $\Delta=-31$, $0.025\cdot 31=0.775$ is past the
cap and $\mu$ drops from $0.35$ to $0.3$---still in $(0,1)$, so Lemma 3.2 applies unchanged to
the only blows the bonus can ever touch (a waiting player's, Section~\ref{sec:policy}).

\subsection{Geometry from the invariants}

\begin{lemreach}[Reach]
Assume (I1)--(I4), and let a position be reached in which every $E_S$ with $S\ni e$ is alive and
$R_e$ holds exactly one player stack, namely slot $e$'s. Then that stack can strike exactly the
enemies $E_S$ with $S\ni e$, and for each such $S$ its only legal approach hex is the docking
$d_S^e$. The starting position is the special case.
\end{lemreach}

\noindent\emph{Proof.} Every hex outside $R_e$ that is adjacent to $R_e$ is either impassable or
holds a living enemy, and neither is enterable (R3, R4): by (I2) the regions are the connected
components of the free hexes once the enemy hexes are removed, so the only free non-region
neighbours of $R_e$ are enemy hexes, and by (I3) those are exactly the $z_S$ with $S\ni e$, alive
by hypothesis. The movement BFS (R11) from the stack's hex therefore cannot leave
$R_e$. By hypothesis $R_e$ holds that one stack alone, so no ally blocks the walk, and
by (I4) the whole of $R_e$ is within movement range. \texttt{WALK\_AND\_ATTACK} needs a free hex
adjacent to the target; by (I3) the only hexes of $R_e$ adjacent to an enemy are the $d_S^e$, and
$d_S^e$ is adjacent to $z_S$ alone. \qedbox

\begin{lemalt}[Alternating triples]
The only pairwise non-adjacent 3-subsets of a hex's six neighbours are the two alternating
triples.
\end{lemalt}

\noindent\emph{Proof.} Consecutive neighbours on the six-cycle are adjacent to each other, in
both row parities; a 3-subset of a 6-cycle with no two consecutive members is an alternating
triple. \qedbox\ (Machine-checked for both parities, \src{scripts/test\_obstacles.py}.)

\smallskip
So (I1) is a constraint, not a design choice: an enemy that three mutually sealed regions must
reach simultaneously is forced to present an alternating triple. This is the one place where hex
geometry does real work, and why the construction would not transfer unchanged to a square grid.

\subsection{Gadgets}

\paragraph{Enemy gadget} for $S\in C$, anchored at $(X,Y)$ with $Y$ even: $z_S=(X,Y)$; the three
dockings are the alternating triple $\texttt{TOP\_LEFT}=(X,Y-1)$, $\texttt{RIGHT}=(X+1,Y)$,
$\texttt{BOTTOM\_LEFT}=(X,Y+1)$; the other three neighbours are impassable. Each docking is
adjacent to $z_S$ and to nothing else in the gadget, so a stack that walks in along one arm
cannot slip round to another. (Picture in Section~\ref{sec:thm3}.)

\paragraph{Element gadget} for $e$: a connected region of free hexes with a two-hex stub at its
end carrying $p_e$.

\paragraph{Corridors} connect the dockings to their regions, laid so that no hex of one region is
adjacent to a hex of another region or to an enemy hex other than at its own docking.

\begin{lemadapt}[Adapters]
An orthogonal drawing delivers the three edges of a set-vertex along three of the four axis
directions, and which three is not ours to choose; the three dockings, by Lemma D.2, are fixed.
For each of the four possible triples of incoming directions there is a $9\times 9$ pattern in
which three pairwise non-touching paths run from the three used boundary midpoints (``ports'') to
the three dockings, each path meeting $z_S$ only at its final hex, with the unused port left
impassable.
\end{lemadapt}

\noindent\emph{Proof.} Exhibited, one pattern per case, in Appendix~\ref{app:adapters}, and
machine-checked hex by hex (\src{scripts/verify\_x3c.py::check\_enemy\_adapters}): each arm is a
path, starts at its port, ends at a distinct member of the alternating triple, touches the enemy
only at its last hex, no two arms share or touch a hex, and the unused port is untouched. The
arms are matched to the dockings in cyclic order, which is what keeps the routing planar inside
the box. \qedbox

\smallskip
The box is $9\times 9$ with the enemy at its centre $(4,4)$ in a local frame whose row 4 is even;
the ports are the four boundary midpoints. The parity of the local frame is discharged in step 5
of Lemma D.4.

\subsection{The embedding}

\begin{lemembed}[Embedding]
Let $(X,C)$ be a \PX{} instance with $N=|X|+|C|$. In time polynomial in $N$ one can compute
either
\begin{itemize}
\item a hex board of polynomial size, a set of impassable hexes, enemy hexes $z_S$ ($S\in C$) and
  deployment hexes $p_e$ ($e\in X$) satisfying (I1)--(I4), or
\item a fixed no-instance $G_{\mathrm{no}}$ of \ALLOC{} (a $1\times 1$ board, one slot, one
  creature of stock one, no enemies, $W=1$), emitted \textbf{only} when the given encoding is
  malformed or fails one of the polynomial-time no-certificates of step 0---$|X|$ not divisible
  by 3, $|C|<q$, an element lying in no set, or a failed planarity test.
\end{itemize}
\end{lemembed}

\noindent ``Only when'', not ``exactly when'': every input sent to $G_{\mathrm{no}}$ is
certifiably a no-instance, but most no-instances are \emph{not} sent there---they receive an
ordinary board whose game answer is no, which is what a many-one reduction should do. (An
``exactly when'' would be a decision procedure for \PX.) The second branch is what makes the map
a \textbf{total} Karp reduction: it outputs a target instance on every source encoding.

\smallskip
\noindent\emph{Proof, in seven steps.} Step 0 disposes of degenerate and malformed inputs;
steps 1--2 are graph surgery; step 3 quotes the drawing literature; and steps 4--6 are the
hex-level construction, where the adjacency bookkeeping lives.

\paragraph{Step 0: degenerate inputs.} Delete repeated members of $C$ (a duplicate 3-set is never
needed by an exact cover, so the answer is unchanged); afterwards the incidence graph $G(X,C)$ is
simple. If $|X|=0$, the instance is a \textbf{yes}-instance (the empty cover works, $q=0$) and
the algorithm outputs the fixed board built for $X=\{1,2,3\}$, $C=\{\{1,2,3\}\}$---a yes-instance
of the game---rather than sending the empty encoding through the geometric steps, whose
bookkeeping assumes at least one element.
If $|X|$ is not divisible by 3, if $|C|<q$, or if some element lies in no member of $C$,
no exact cover exists: output $G_{\mathrm{no}}$. Otherwise every set-vertex of $G$ has degree
exactly 3 and every element vertex degree $d_e\ge 1$. Compute a combinatorial planar embedding of
$G$---a rotation system with a choice of outer face---in linear time (Hopcroft--Tarjan~\cite{HT74}, or
Boyer--Myrvold~\cite{BM04}, which returns the embedding directly). If $G$ is not planar, the encoding is not
a \PX{} instance and maps to $G_{\mathrm{no}}$, which keeps the reduction total under either
convention for defining the promise.

\paragraph{Step 1: degree reduction preserving the rotation.} Replace every element vertex $e$ of
degree $d$ by a path $v_e^1 - \cdots - v_e^d$ of $d$ new vertices, attaching the $i$-th incident
edge of $e$ (in the cyclic order given by the embedding, cut at an arbitrary point) to $v_e^i$.
Call the result $G'$. \emph{Claim: $G'$ is plane, $\Delta(G')\le 3$, and an embedding extending
the given one outside a small disc around each $e$ is computable in linear time.} Degrees: a
set-vertex keeps degree 3; an interior path vertex has two path neighbours and one incidence
edge; an endpoint has degree $\le 2$. Embedding: fix a closed disc $\Delta_e$ around $e$ meeting
no other vertex and meeting the incident edges in $d$ initial segments; the boundary circle meets
them at points $b_1,\dots,b_d$ in the rotation's cyclic order. Cut the circle open between $b_d$
and $b_1$ and straighten the disc to a rectangle whose top side carries $b_1,\dots,b_d$ left to
right; draw the path horizontally across the middle with $v_e^i$ below $b_i$ and join each
$v_e^i$ to $b_i$ vertically. The vertical segments are pairwise disjoint and meet the horizontal
one only at their own endpoints, so the picture inside the rectangle is plane; outside $\Delta_e$
nothing changed. The \emph{cyclic} order of the incidence edges around $e$ becomes the
\emph{linear} order along the path, and the cut is made in a face corner, so no crossing is
created. Sizes: $|V(G')|=4|C|$ and $|E(G')|=6|C|-|X|$, both $O(N)$.

\paragraph{Step 2: connected components.} $G'$ need not be connected (the incidence graph of
$X=\{1,\dots,6\}$, $C=\{\{1,2,3\},\{4,5,6\}\}$ is a planar yes-instance with two components).
Let $G'_1,\dots,G'_t$ be its components, $\sum_i n_i=4|C|$; run steps 3--6 on each separately,
obtaining boards $B_i$ of polynomial size, then \textbf{pack}: place them left to right,
separated by vertical strips of impassable hexes two hexes wide and surrounded by an impassable
border two hexes thick, with each block's vertical offset chosen even (the border allows it).
Hexes of distinct blocks then differ by at least 2 in a coordinate, and two hexes at $L_\infty$
distance $\ge 2$ are never adjacent (each hex neighbour differs by at most 1 in each coordinate,
R1), so no invariant can be violated across blocks and each region stays inside its own block.
The packed board remains polynomial. (The implementation packs one level earlier:
\src{planar\_embed.py} places the component \emph{drawings} side by side with two empty grid
columns between them and scales once, so features of distinct components end at $L_\infty$
distance at least $2\lambda-2\rho=32\ge 2$ by the step-4 separation inequality, and the hexes
between them---impassable by default---form exactly the strip this step prescribes. This is the
third named subroutine substitution of D.6; the (SEP$'$) check measures every non-incident pair
of features, cross-component pairs included, so the substituted packing is verified directly on
every built board.)

\paragraph{Step 3: an orthogonal grid drawing.} For a connected plane graph $H$ with
$\Delta(H)\le 4$ on $n$ vertices we need a drawing with \textbf{(D1)} vertices at distinct points
of $\mathbb{Z}^2$; \textbf{(D2)} edges as rectilinear grid paths leaving each endpoint along one
of the four axis directions, each (vertex, direction) pair used at most once; \textbf{(D3)} two
edge paths meeting only at a shared endpoint and no path passing through a vertex; \textbf{(D4)}
all coordinates polynomial in $n$; \textbf{(D5)} computable in polynomial time. Nothing
else---not bend-minimality, not area optimality, not preservation of the step-1 embedding (the
correctness argument never refers to it again). Tamassia--Tollis~\cite{TT89}, in the form of
\cite[Thm.~7.3]{DG13}, supplies exactly this for \textbf{connected} 4-plane graphs---$O(n^2)$
area, at most four bends per edge, linear time---and step 2 has already reduced us to the
connected case, with $\Delta(G')\le 3\le 4$. The cited theorem bounds the \textbf{area} by
$O(n^2)$; it does not bound either side by $O(n)$, and we do not need it to. Enclose the drawing
$\Gamma$ in its bounding box and write $g$ for the larger side, so $g=O(n^2)$ in the worst case;
every claim below uses only that $g$ is polynomial, so the board of $O(g)\times O(g)$ hexes is
polynomial and the unary speed $\sigma$ is too. Strong NP-hardness is unaffected. (An earlier
draft asserted a $g\times g$ grid with $g=O(n)$, which the cited theorem does not give; the
polynomial bound is all the argument uses.)

\paragraph{Step 4: scaling, and the separation inequality.} Let $\lambda:=20$ (the scale factor,
even) and $\rho:=4$ (the gadget-box radius), and map the drawing into hex coordinates by
\[
  \Phi(i,j):=\bigl(\lambda\cdot(i-i_{\min})+\omega,\ \lambda\cdot(j-j_{\min})+\omega\bigr),
  \qquad \omega:=6,
\]
where $(i_{\min},j_{\min})$ is the lower-left corner of $\Gamma$'s bounding box; $\omega$ is a
coordinate offset and has nothing to do with the damage multiplier $\mu$ of
Section~\ref{sec:thm3}. \textbf{The offset $\omega$ must exceed the box radius $\rho$}: an
earlier draft used $\omega=2$, under which a legal drawing vertex at the bounding-box corner maps
to $(2,2)$ and its $9\times 9$ box reaches $(-2,-2)$, off the board---a genuine construction
failure, not a cosmetic one. With $\omega=6$ every box lies at coordinates $\ge 2$, and $\omega$
even keeps every image row $\lambda(j-j_{\min})+\omega$ even. Fill in each drawn edge: a unit segment of $\Gamma$ becomes the $\lambda+1$ hexes on the
corresponding axis-aligned run, which is legitimate because the 4-neighbour square grid is a
subgraph of hex adjacency in both row parities---$(x,y)$ is adjacent to $(x\pm 1,y)$ and to
$(x,y\pm 1)$ in either parity (R1; machine-checked in \src{test\_obstacles.py}).

Two features of $\Gamma$ that share no point are at $L_\infty$ distance $\ge 1$ (their segments
lie on integer grid lines with integer endpoints), and $\Phi$ scales distances by exactly
$\lambda$, so disjoint features land at $L_\infty$ distance $\ge\lambda=20$ \textbf{(SEP)}.
Step 5 replaces a $(2\rho+1)\times(2\rho+1)$ box around each vertex image by a gadget and
truncates every corridor at the boxes at its ends; every hex of a box is within $L_\infty$
distance $\rho$ of its vertex image, so for two \emph{non-incident} features---two boxes of
distinct vertices, a box and a corridor of an edge not incident to it, two corridors of edges
sharing no endpoint---
\begin{equation*}
  L_\infty\text{-dist}\ \ge\ \lambda-2\rho\ =\ 20-8\ =\ 12\ \ge\ 2,
  \tag{SEP$'$}
\end{equation*}
and two hexes at $L_\infty$ distance $\ge 2$ are never adjacent. (An earlier draft used
$\lambda=9$: $9-8=1$, and two boundary cells of unrelated boxes could touch---the repair is this
inequality.) Two corridors \emph{incident to the same box} leave it through different ports (by
(D2)); outside the box each runs straight along its own axis for at least $\lambda-\rho=16$ hexes
before its first possible bend (bends of $\Gamma$ sit at grid points, whose images are
$\ge\lambda$ away from the vertex image); a point at parameter $t$ on one run and $t'$ on the
other are at $L_\infty$ distance $\max(t,t')$ (perpendicular axes) or $t+t'$ (opposite axes), and
outside the box $t,t'\ge\rho+1=5$, so in either case the distance is $\ge 5\ge 2$ and the runs
never touch. (Both
cases occur: every three of the four axis directions contain exactly one opposite pair.)

\paragraph{Step 5: the gadget boxes.} \emph{Set-vertex $S$.} Its image sits on an even row.
Replace the $9\times 9$ box around it by the Lemma D.3 adapter for the triple of axis directions
along which $\Gamma$ delivers the three incident edges; the centre becomes $z_S$, the arms are
assigned to the elements of $S$ by which port each element's corridor arrives at, and every other
hex of the box, the unused port included, is impassable. \emph{Parity:} the adapter is stated in
a local frame with row 4 even; the translation onto the board shifts rows by $Y-4$ with $Y$ even,
and hex adjacency in offset coordinates depends only on row parity and is invariant under column
translation (R1), so the pattern carries over adjacency for adjacency---this is why $\lambda$ and
$\omega$ are even.

\emph{Element-path vertex $v_e^i$.} Its degree in $G'$ is at most 3, so at least one axis
direction is unused. The free hexes of its box are the axis segments from the centre to the used
ports plus the centre: a ``plus'' with at most three arms, connected, contained in the box,
meeting the boundary exactly at the used ports. Adjacencies inside the plus are
irrelevant---every hex of it belongs to the single region $R_e$---which is exactly why the
set-vertex needs a hand-built adapter and the element vertex does not.

\emph{The deployment stub.} Take $v_e^1$ and one axis direction $u$ unused at it (one exists:
$\deg(v_e^1)\le 2$). Declare free the two hexes at distance 1 and 2 from the centre along $u$,
and set $p_e$ to the one at distance 2. Then $p_e$ is connected to the plus, so it lies in $R_e$;
and $p_e$ has \textbf{exactly one} free neighbour. For $u=\texttt{RIGHT}$: $p_e=(X+2,Y)$ with $Y$
even has neighbours $(X+1,Y)$ (the stub's other hex, free), $(X+3,Y)$ (beyond the two-hex stub,
impassable), and four hexes in rows $Y\pm 1$, where the only free hexes of the box are the
vertical-axis cells $(X,Y\pm 1)$, each at $L_\infty$ distance $\ge 2$ from all four. For a
vertical $u$, say $p_e=(X,Y+2)$: rows $Y+1$ and $Y+2$ contain no free hex except the stub itself,
and $(X,Y+3)$, $(X+1,Y+3)$ lie beyond the stub. The cases $u=\texttt{LEFT}$ and $u=\texttt{UP}$
are checked the same way, hex by hex; the column case is \emph{not} a reflection of the row
case---reflecting in a column does not preserve hex adjacency in offset coordinates. The clause
``exactly one free neighbour'' holds for every used-port set and every unused $u$, in both row
parities: the finitely many local configurations are enumerated in the apparatus
(\src{check\_stub.py}, adjacency taken from the model's own neighbour rule, 64 cases, 0
failures). Every free hex within $L_\infty$ distance 12 of the stub belongs to a feature
incident to $v_e^1$, hence to $R_e$: the corridors leaving $v_e^1$ through its used ports
come as close as $L_\infty$ distance 5 on the instance corpus---measured by
\src{verify\_embedding.py}, which also checks that membership on every corpus board---and,
lying outside the box of $v_e^1$, they stay at $L_\infty$ distance $\ge 3$ from a stub two
hexes from its centre; every non-incident feature is at distance $\ge 12$ by (SEP$'$). So the
stub creates no adjacency anywhere else.

\paragraph{Step 6: closing the board.} Declare impassable every hex not made free in steps 4--5;
set $\sigma$ to the hex count of the packed board (legal: creature statistics are input,
Section~\ref{sec:generalized}). The invariants hold: \textbf{(I1)} the free neighbours of $z_S$
are exactly the three arm ends---the alternating triple, pairwise non-adjacent by
Lemma D.2---and no hex outside the box is adjacent to $z_S$ by (SEP$'$). \textbf{(I2)} the pluses
of $e$'s path vertices, joined by the path-edge corridors, with the incidence corridors, the
adapter arms ending at $e$'s dockings and the stub, form one connected set; distinct regions never
touch---their features are non-incident in $\Gamma$, so (SEP$'$) applies, and inside an adapter
box Lemma D.3 has checked the three arms pairwise non-touching; every free non-enemy hex belongs
to exactly one element, so there are exactly $3q$ components. \textbf{(I3)} $R_e$ contains $p_e$;
its hexes adjacent to an enemy are exactly the dockings $d_S^e$ for $S\ni e$ (inside $S$'s box
only the last hex of each arm touches $z_S$; outside, (SEP$'$)).
\textbf{(I4)} $\sigma$ is the hex count, which is $\ge|R_e|$ for every $e$. \qedbox

\subsection{Correctness}

Throughout, fix a feasible allocation $c:X\to\mathbb{Z}_{\ge 0}$ with $\sum_e c_e\le 3q$ and an
arbitrary play of round 1 against $(\ddagger)$.

\begin{lemnoint}[No interference]
No player creature dies during round 1, and every player stack gets its action.
\end{lemnoint}

\noindent\emph{Proof.} Under $(\ddagger)$ the defence never initiates. An enemy retaliates only
when struck, at most once per round (R12; neither \texttt{WAIT} nor \texttt{DEFEND} consumes the
charge, Section~\ref{sec:policy}), and its strike deals $\max(1,\lfloor 1\cdot 1\cdot
1\cdot 1\rfloor)=1$ damage, since $\mathit{att}(Q)=\mathit{def}(P)=1$ gives $\Delta=0$. A player stack of
$c\ge 1$ creatures of 4 hit points absorbs 1 damage without losing a creature (R5, R6). Each
player stack strikes at most once, so it receives at most one retaliation. Every living stack
takes its terminal action once per round, whichever phase it lands in (R9). \qedbox

\begin{lembudget}[Budget]
Let the play kill $t$ enemies. Then $t\le q$. If $t=q$, then $\sum_e c_e=3q$; every slot holds
exactly one creature; every deployed stack strikes an enemy that dies; and each dead enemy is
struck by exactly three distinct one-creature stacks.
\end{lembudget}

\noindent\emph{Proof.} Each enemy is a single creature of 3 hit points, so it dies exactly when 3
damage has accumulated on it (R8), and that damage comes only from the stacks that struck it. A
stack \textbf{strikes} at most once per round---\texttt{WALK\_AND\_ATTACK} ends its turn, and
\texttt{WAIT} defers the turn without granting a second strike (R9)---so the striker sets of
distinct enemies are disjoint. (The weaker phrasing ``a stack acts once per round'' would be
false under \texttt{WAIT}; the argument runs on strikes throughout.) Let $K$ be the set of dead
enemies, $|K|=t$, and $A_S$ the strikers of $S\in K$. By the one-round lemma of
Section~\ref{sec:policy} the blow of a stack of $c$ creatures delivers at most $D(c)$, so
Lemma 3.2 applied to $(c_a)_{a\in A_S}$ gives $\sum_{a\in A_S}c_a\ge 3$. Summing over the
disjoint $A_S$:
\begin{equation*}
  3t\ \le\ \sum_{S\in K}\sum_{a\in A_S}c_a\ \le\ \sum_e c_e\ \le\ 3q,
  \tag{$\dagger$}
\end{equation*}
whence $t\le q$. Suppose $t=q$, so every inequality in $(\dagger)$ is tight. The middle one tight
means every allocated creature sits in a stack that struck a member of $K$---the left sum counts
exactly the creatures in striker stacks of $K$, the right one all allocated creatures. The first
one tight, together with $\sum_{a\in A_S}c_a\ge 3$ for each $S$, forces $\sum_{a\in A_S}c_a=3$
for every $S$; the equality case of Lemma 3.2 then gives $|A_S|=3$ and $c_a=1$ for each. Hence
the $3q$ allocated creatures sit in exactly $3q$ singleton stacks; a stack occupies one slot and
there are exactly $3q=k$ slots, so every slot holds exactly one creature and each struck a member
of $K$. \qedbox

\smallskip
Lemma D.6 uses no geometry whatsoever.

\begin{lemconf}[Confinement]
In the situation of Lemma D.6 with $t=q$, when the stack of slot $e$ takes its terminal action,
every enemy $E_S$ with $S\ni e$ is still alive. Consequently, by Lemma D.1, that stack can strike
only enemies $E_S$ with $S\ni e$.
\end{lemconf}

\noindent\emph{Proof.} Induction on the realized order of \textbf{terminal actions}---the single
move, attack, or defend each stack performs. That order is fixed by the phase structure (R9),
with a stack that issues \texttt{WAIT} taking its terminal action at its later position in the
\texttt{WAIT} phase; the argument uses only that each stack takes exactly one terminal action, at
one position in that order.

Suppose the claim holds for every stack that took its terminal action before slot $e$. Then every
strike so far was delivered from a docking of the striker's own region: by (I3) the only hexes of
$R_a$ adjacent to an enemy are $a$'s own dockings, and by the inductive hypothesis every earlier
striker was still confined to its region when it moved, because all enemies bounding that region
were alive at that moment.

Let $E_S$ be any enemy already dead, and suppose $e\in S$. By Lemma D.6 every stack in play holds
one creature and delivers at most $D(1)=1$ per blow; $E_S$ has 3 hit points and is struck by
exactly three stacks over the round, so each blow delivered exactly 1 and $E_S$ died on the
third---all three strikers have already acted. Each struck from a hex adjacent to $z_S$; by
Lemma D.5 no player stack dies, and by R11 a stack that performed \texttt{WALK\_AND\_ATTACK}
remains on its approach hex for the rest of the round, so no two strikers can have used the same
hex; hence those hexes are the three free neighbours of $z_S$: the dockings
$d_S^{e'}$, one per $e'\in S$. By the previous paragraph the stack that struck from $d_S^{e'}$ is
the stack of slot $e'$. In particular the stack of slot $e$ is among them, so it has already
taken its terminal action---contradicting that it is taking it now. Hence no $E_S$ with $S\ni e$
is dead, the boundary of $R_e$ is intact, and Lemma D.1 applies unchanged. \qedbox

\smallskip
Why this lemma is needed at all: a dead unit stops blocking its hex (R4), so the regions are
\emph{not} permanently separated---every kill opens a doorway between three of them. Lemma D.7
says the doorway only ever opens for stacks that have already spent their terminal action;
without it the no-direction would leak. The induction covers plays in which a stack moves without
attacking, defends, or waits---and the three cases are not the same: a stack that moves or
defends has \emph{spent} its terminal action, while a stack that waits has \emph{postponed} it,
and the induction simply reaches it at its later position. The argument never assumes an acting
stack strikes, and the one-round lemma of Section~\ref{sec:policy} covers the only other effect
of waiting---a possible damage reduction against a defended target, which only helps the
no-direction.

Lemma D.5 enters the proof only through ``two stacks cannot occupy one hex'', and the argument
survives without it: count strikers instead of hexes. The three blows on $E_S$ came from three
distinct stacks (each strikes at most once, Lemma D.6), each from a docking of its own region
by the first paragraph, and the dockings adjacent to $z_S$ are $d_S^{e'}$ for $e'\in S$, one
per region---so the strikers are the stacks of the three slots of $S$, and slot $e$ is among
them. We keep Lemma D.5 because it is also what makes the phase-by-phase picture of
Section~\ref{sec:thm3} literal: nobody dies, nobody vacates a docking.

\begin{lemyes}[yes $\Longrightarrow$ yes]
If $(X,C)$ has an exact cover, $G_3(X,C)$ is a yes-instance.
\end{lemyes}

\noindent\emph{Proof.} Let $C'=\{S_1,\dots,S_q\}$ be an exact cover. Allocate one creature to
every slot (total $3q$, exactly the stock). Each $e$ lies in exactly one $S\in C'$; let slot $e$
issue \texttt{WALK\_AND\_ATTACK} to $d_S^e$ against $E_S$. No stack waits, so every blow lands in
the \texttt{NORMAL} phase, meets the un-raised defence (Section~\ref{sec:policy}) and deals its
nominal $D(1)=1$. Every move is legal whatever order the engine imposes: only the three slots of
$S$ ever target $E_S$, so when slot $e$ acts, $E_S$ has absorbed at most 2 and is alive; and
$d_S^e$ is free, because in this play every stack performs exactly one
\texttt{WALK\_AND\_ATTACK}, ending on a docking of its own region, and $R_e$ holds only slot
$e$'s stack (I2, I3), so doorways opened by earlier kills are never used. Each $S\in C'$ accumulates 3 damage and dies; the destroyed value is $q=W$.
\qedbox

\begin{lemcover}[yes $\Longrightarrow$ exact cover]
If $G_3(X,C)$ is a yes-instance, $(X,C)$ has an exact cover.
\end{lemcover}

\noindent\emph{Proof.} Take an allocation and a play destroying value $\ge q$. Every enemy has
value 1, so the number $t$ of dead enemies is $\ge q$; by Lemma D.6, $t\le q$, hence $t=q$ and
the tightness half of Lemma D.6 applies. Let $K$ be the set of dead enemies. Each $S\in K$ is
struck by exactly three one-creature stacks, and by Lemma D.7 each striker of $E_S$ is the stack
of a slot $e\in S$; the three strikers are distinct and $|S|=3$, so they are exactly the slots of
$S$. A slot strikes once, so distinct members of $K$ use disjoint slot sets; $|K|=q$ triples use
all $3q$ slots. Hence $K$ is a family of $q$ pairwise disjoint 3-sets covering $X$. \qedbox

\smallskip
{\sloppy
\noindent\textbf{Proof of Theorem 3.} Lemmas D.8 and D.9 give the equivalence; Lemma D.4 gives a
total polynomial-time construction; every number in $G_3(X,C)$ is polynomially bounded, so the
reduction is polynomial even under unary encoding, and \PX{} is NP-complete~\cite{DF86}---a
problem with no numeric parameters, so the hardness obtained through it is strong.
Membership is Lemma 2.1. \qedbox\par}

\smallskip
\noindent\textbf{Proof of Corollary 3.1.} Fix the allocation $c_e=1$ for every $e$, so the
instance has no allocation decision and the only choice left is the play. The yes-direction is
Lemma D.8, which uses precisely this allocation. For the no-direction, with $c\equiv 1$ every blow
delivers at most $D(1)=1$; an enemy has 3 hit points, so each kill consumes strikes of at least
three distinct stacks, and striker sets of distinct dead enemies are disjoint, so killing $t$
enemies consumes at least $3t$ of the $3q$ stacks and $t\le q$; at $t=q$ every dead enemy is
struck by exactly three singleton stacks and every stack strikes a dead enemy---the conclusion of
Lemma D.6, obtained here without the resource lemma's equality case (which is only needed to rule
out unequal stack sizes). Lemmas D.7 and D.9 then apply verbatim. \qedbox

\subsection{Scope of the machine checks}\label{app:mcscope}

Lemma D.2, Lemma D.3's four adapters, and the square-grid-inside-hex-adjacency fact of step 4 are
machine-checked directly. The embedding algorithm of Lemma D.4 is implemented step by step with
the lemma's literal constants ($\lambda=20$, $\rho=4$, the D.3 adapters unchanged):
\src{embed\_lemma.py} is steps 0--6, on top of the planarity and drawing machinery of
\src{planar\_embed.py}---with three subroutine substitutions, named here so that a reader who
opens the artifact is not surprised. The planarity test is Demoucron--Malgrange--Pertuiset~\cite{DMP64} (cubic,
chosen because it yields the face set directly), not the linear-time Hopcroft--Tarjan or
Boyer--Myrvold that step 0 cites; the drawing is a from-scratch st-ordered visibility
construction for maximum degree 3, not \cite{TT89} as quoted in step 3; and step 2's packing of
finished boards is realized one level earlier, as a packing of the component drawings (two empty
grid columns between components, $2\lambda$ hexes after scaling---sound by the step-4 separation
inequality, and measured directly by (SEP$'$), which checks every non-incident pair of features,
cross-component pairs included). All three substitutions are
legitimate because the proof consumes only (D1)--(D5), polynomial time and the separation
inequality, and \src{validate\_drawing} machine-checks the drawing properties on every build.
\src{verify\_embedding.py} feeds every instance family of the published corpora through the
algorithm: every board built satisfies (I1)--(I4) and the feature-based (SEP$'$) separation
check, every
degeneracy skip carries a certificate re-verified against the instance (non-planarity is the DMP
test's own verdict, cross-checked by a planted non-planar control), the $G_{\mathrm{no}}$ and
$|X|=0$ shortcut branches are exercised---$G_{\mathrm{no}}$ is itself played out as a genuine
no, and a battery of malformed encodings (wrong arity, non-integers, out-of-range or repeated
members, a non-sequence, negative sizes) must each produce a certified $G_{\mathrm{no}}$ rather
than a crash---and on the smallest boards
the full game search runs end-to-end---under the historical and the published constants---and
agrees with \textsc{X3C}. The big
exhaustive-search suites still take their boards from a compact router (placement by
hill-climbing, corridors by BFS with clearance), because the lemma's literal boards are $\lambda$
times larger per drawing unit than the router's; the router reports failure rather than emitting
a board that violates the invariants. What remains a hand proof is the lemma's universal
claim---that the algorithm succeeds on \emph{every} planar instance---the scope note of
Section~\ref{sec:contrib}, item~7. The correctness suite additionally verifies, on every built yes-instance,
that the unique winning allocation is the all-ones vector predicted by Lemma D.6, and runs a
negative control with $\mathit{def}(Q)=\mathit{att}(P)$---destroying Lemma 3.2---under which a
no-instance of \textsc{X3C} does turn into a yes-instance of the game, so a pass carries evidence
rather than silence.

\section{Full proofs of Theorems 1, 2 and 4}\label{app:thm124}

This appendix makes Theorems 1, 2 and 4, Proposition 1.1 and Corollaries 4.1 and 4.2
self-contained, in the way Appendix~\ref{app:thm3} does for Theorem 3: the complete
constructions with every creature tuple written out, the geometry, the damage accounting,
and both directions of each equivalence. The statements are repeated here so that the
appendix can be read on its own; the model is Section~\ref{sec:model}, the policy
$(\ddagger)$ and the one-round lemma are Section~\ref{sec:policy}, and rule tags
\texttt{R1}--\texttt{R17} resolve to engine lines in Appendix~\ref{app:rules}.

Three of the repairs below change what the body says rather than only how it says it, and
they are flagged where they occur: Theorem 1's no-direction is re-derived over disjoint
striker sets and uses no reach hypothesis at all (the Theorem 1 section); Proposition
1.1's matching hypothesis is quantified over \emph{every} position reachable in round-1
play, not only over the starting position, because a dead unit stops blocking its hex
(the Proposition 1.1 section); and all three reductions are made total (the conventions
section).

\subsection*{Conventions common to the three constructions}

\paragraph{The grid.} We use the offset (``even-row shifted'') coordinates of R1
throughout, in the concrete form printed in Appendix~\ref{app:adapters}: writing $(x,y)$
for column $x$ of row $y$, the six neighbours of $(x,y)$ are
\[
  (x\pm 1,\,y),\quad (x-\epsilon,\,y-1),\ (x+1-\epsilon,\,y-1),\quad
  (x-\epsilon,\,y+1),\ (x+1-\epsilon,\,y+1),
  \qquad \epsilon:=y\bmod 2 .
\]
So from an \textbf{even} row the two upper neighbours sit at columns $x$ and $x+1$, and
from an \textbf{odd} row at columns $x-1$ and $x$. Distance is the axial formula of R2:
with $A(x,y):=x+\lfloor y/2\rfloor$, $\delta A:=A(q)-A(p)$ and $\delta y:=y_q-y_p$,
\begin{equation*}
  \mathrm{dist}(p,q)=
  \begin{cases}
    \max(|\delta A|,|\delta y|) & \text{if } \delta A\ \text{and}\ \delta y\
      \text{are both}\ \ge 0\ \text{or both}\ <0,\\[2pt]
    |\delta A|+|\delta y| & \text{otherwise.}
  \end{cases}
  \tag{R2}
\end{equation*}
Every coordinate computation below is done in $(A,y)$.

\paragraph{$(\star)$ in full.} $(\star)$ is the specialization announced in
Section~\ref{sec:results}, in its complete form: fix $\alpha:=1$ and give
\emph{every player type} $\mathit{att}=\mathit{def}=1$, \emph{every enemy type}
$\mathit{att}=\mathit{def}=1$, and every enemy type flat damage
$\mathit{dmg}_{\min}=\mathit{dmg}_{\max}=1$. Then $\Delta=0$ on a player$\to$enemy blow
\emph{and} on the retaliation, so $f_{\mathrm{att}}=f_{\mathrm{def}}=1$ in both directions
(R7) and a stack of $c\ge 1$ creatures of flat per-creature damage $d\ge 1$ delivers
nominal damage
\begin{equation*}
  \mathrm{dmg}=\max\bigl(1,\lfloor c\,d\rfloor\bigr)=c\,d .
  \tag{$\star$}
\end{equation*}
The lower bound $d\ge 1$ is not decoration: it is the model's own domain restriction
$1\le\mathit{dmg}_{\min}\le\mathit{dmg}_{\max}$ (Section~\ref{sec:rules}), and it is what
makes the clamp $\max(1,\cdot)$ inert here. Every claim below that says ``nominal is
dealt'' uses $c\ge 1$ and $d\ge 1$; every claim that quantifies over an arbitrary play says
``at most nominal'', as Section~\ref{sec:policy} requires.

An earlier form of $(\star)$ fixed only the player's attack and the enemy's defence, which
left three components of the constructed tuples undetermined---the player's defence and the
enemies' attack and damage---so that the reductions did not output a single well-defined
\ALLOC{} instance. The form above fixes them, and Lemma E.2 is where they are used.

\paragraph{Complete tuples.} A creature type is the tuple
$(\mathit{att},\mathit{def},\mathit{dmg}_{\min},\mathit{dmg}_{\max},\mathit{hp},
\mathit{spd},\mathit{flags})$ together with its \emph{value} (Section~\ref{sec:rules}).
In all three constructions $\mathit{flags}=\emptyset$: every creature is melee, single-hex,
non-shooting, with the default single retaliation charge---the fragment
\texttt{H3-det-melee} of Section~\ref{sec:rules}. Every player type has $\mathit{value}=0$,
so the objective reads enemy values only, as the Problem of Section~\ref{sec:problem}
prescribes. With $(\star)$ fixing four of the numbers and $\mathit{flags}$ and
$\mathit{value}$ fixed here, each construction below has only $d$, $\mathit{hp}$,
$\mathit{spd}$ and the stock left to name, and it names all four.

\paragraph{Totality.} A Karp reduction must output a target instance on every source
encoding, and the constructions below presuppose a nonempty, well-formed source instance.
All three are made total by the same three-way branch, in the pattern of Lemma D.4:

\begin{enumerate}
\item If the encoding is malformed, or fails one of the syntactic checks listed with the
  construction---each of which certifies in polynomial time that no solution exists---output
  the fixed no-instance $G_{\mathrm{no}}$ of Lemma D.4: a $1\times 1$ board, one slot, one
  creature of stock one, no enemies, $W=1$. Its destroyed value is $0<1$, so it is a
  no-instance, and $W=1\in\mathbb{Z}_{>0}$ as the Problem requires.
\item If the source instance is \emph{empty}---no numbers at all---it is a
  \textbf{yes}-instance (the empty selection works), so output instead the instance the
  construction builds from a fixed nonempty source yes-instance, named with each
  construction. Sending the empty encoding through the geometric steps is what must be
  avoided: it yields an empty board and $W=0\notin\mathbb{Z}_{>0}$.
\item Otherwise, run the construction.
\end{enumerate}

\noindent``Only when'', not ``exactly when'', again: every encoding routed to
$G_{\mathrm{no}}$ is certifiably a no, and most no-instances are not routed there but
receive an ordinary board whose game answer is no.

\begin{lemEblow}[one blow, and where it is issued]
Fix any allocation and any play of round 1 against $(\ddagger)$ in any of the three
constructions. Then every player stack delivers \textbf{at most one blow} in the round, and
if it delivers one, it is issued at the stack's terminal action, from the stack's own
deployment hex.
\end{lemEblow}

\noindent\emph{Proof.} A stack has exactly one terminal action per round---move, attack, or
defend---and \texttt{WAIT} is not one: it postpones the terminal action into the
\texttt{WAIT} phase without moving or striking (R9, Section~\ref{sec:policy}). The only
player action that strikes is \texttt{WALK\_AND\_ATTACK}, which moves and strikes as a
single terminal action (R11), so at the moment it is issued the stack still stands where it
stood at the start of the round, namely on its deployment hex. The only other blow a player
stack could deliver is a retaliation, and under $(\ddagger)$ the defence never initiates an
attack (Section~\ref{sec:policy}), so no player stack ever retaliates. Every type used
below has $\mathit{flags}=\emptyset$, so no ability grants a second strike. \qedbox

\smallskip
Lemma E.1 is the step that an earlier draft of Section~\ref{sec:thm4} compressed into ``each type has
stock one, so each stack strikes at most one enemy''. Stock one gives that each \emph{type}
occupies at most one slot; what gives one blow per stack is one terminal action per round
together with the absence of enemy initiative. The distinction is exactly the one whose
neglect produced the first, wrong version of Theorem 1 (Section~\ref{sec:policy}).

\begin{lemEretal}[retaliation is inert]
In round 1 of any of the three constructions, under $(\star)$ and $(\ddagger)$: a player
stack takes at most one retaliation blow; that blow lands strictly after the stack's own
blow; it deals exactly $1$ damage; and no player creature dies.
\end{lemEretal}

\noindent\emph{Proof.} Under $(\ddagger)$ an enemy never initiates, so the only damage a
player stack can take is retaliation, which by R12 resolves \emph{after} the attacker's
blow has been applied. A stack strikes at most once (Lemma E.1), so it draws at most one
retaliation. Under $(\star)$ the retaliating enemy is a single creature of flat damage 1
with $\Delta=\mathit{att}(\text{enemy})-\mathit{def}(\text{player})=1-1=0$, so its blow
deals $\max(1,\lfloor 1\cdot 1\cdot 1\cdot 1\rfloor)=1$ (R7). Every player type below has
$\mathit{hp}=5$, so a stack of $c\ge 1$ creatures absorbs 1 point in its pool
($\mathit{firstHPleft}$ drops from 5 to 4) and its $\mathrm{count}$ is unchanged (R5, R6).
\qedbox

\smallskip
Two consequences are used without further comment. No blow already delivered is affected by
a retaliation, so the damage bookkeeping of each construction may ignore retaliation
entirely; and the objective reads the values of enemy types only
(Section~\ref{sec:problem}), so player losses---of which there are none---could not affect
it in any case.

\begin{lemEmetric}[reach is bounded by the hex metric]
Let a player stack of speed $s$ stand on hex $p$ in some position of the battle, and suppose
it strikes an enemy occupying hex $z$. Then $\mathrm{dist}(p,z)\le s+1$, where
$\mathrm{dist}$ is the R2 distance. The bound holds \textbf{whatever hexes are free at the
time}.
\end{lemEmetric}

\noindent\emph{Proof.} \texttt{WALK\_AND\_ATTACK} walks a path of enterable hexes of length
at most $s$ from $p$ to a hex $h$ adjacent to $z$, and strikes from $h$ (R10, R11). It
suffices that the R2 distance is the graph distance of the R1 adjacency, since a BFS
distance over any subgraph is at least the graph distance over the whole grid; then
$\mathrm{dist}(p,h)\le s$ and $\mathrm{dist}(p,z)\le s+1$.

In the coordinates $(A,y)=(x+\lfloor y/2\rfloor,\,y)$ the six R1 steps become, in \emph{both}
row parities, exactly
\[
  \pm(1,0),\qquad \pm(0,1),\qquad \pm(1,1)
\]
(the two same-row steps give $\pm(1,0)$; the step to column $x+1-\epsilon$ of row $y-1$ and
the step to column $x-\epsilon$ of row $y+1$ give $(0,-1)$ and $(0,+1)$; the remaining two
give $-(1,1)$ and $+(1,1)$). Let $H$ be the hexagon with these six vectors as its vertices.
The R2 formula is the Minkowski gauge of $H$: on the two quadrants where $\delta A$ and
$\delta y$ agree in sign, the boundary of $H$ is the polyline
$\max(|\delta A|,|\delta y|)=1$, and on the other two it is $|\delta A|+|\delta y|=1$. A
gauge is subadditive and takes the value 1 on each of the six steps, so a walk of length
$L$ realizes a displacement of gauge at most $L$. Conversely, when $\delta A$ and $\delta y$
agree in sign, $\min(|\delta A|,|\delta y|)$ diagonal steps with vertical or same-row steps
for the remainder realize the displacement in $\max(|\delta A|,|\delta y|)=\mathrm{dist}$
steps; when they differ in sign no diagonal helps, and $|\delta A|$ same-row with
$|\delta y|$ vertical steps realize it in $|\delta A|+|\delta y|=\mathrm{dist}$ steps.
\qedbox

\smallskip
Lemma E.3 is the tool that disposes of the dynamic-reach objection, and it is worth saying
why it is stated in this reach-independent form. A dead unit stops blocking its hex (R4), so
every kill can open a doorway, and a reach argument that reads off the \emph{starting} free
graph proves nothing about later positions---this is the phenomenon Theorem 3 spends
Lemma D.7 on. In Theorems 1, 2 and 4 no such lemma is needed, because the separation is
metric: the bound of Lemma E.3 does not look at which hexes are free, so freeing hexes
cannot violate it.

\subsection*{Theorem 1}

\begin{thmone}
\ALLOC{} is NP-complete, already for instances with $R=1$, a \textbf{single creature type} in
the player's army, one creature per enemy stack, no obstacles, and a battlefield of one row.
\end{thmone}

\paragraph{Source problem.} \PARTITION{} (\cite[SP12]{GJ79}, \cite{Kar72}): given positive
integers $a_1,\dots,a_n$ with $\sum_i a_i=2B$, is there $S\subseteq[n]$ with
$\sum_{i\in S}a_i=B$?

\paragraph{The instance $G(a)$.} Given $a=(a_1,\dots,a_n)$ with $n\ge 1$, every
$a_i\in\mathbb{Z}_{>0}$ and $\sum_i a_i=2B$:

\begin{itemize}
\item \textbf{Battlefield.} One row of $5n$ hexes, indexed $0,\dots,5n-1$; no obstacles.
  Block $j$ ($1\le j\le n$) occupies hexes $5(j-1),\dots,5j-1$, and within it
  \[
    p_j:=5(j-1)\quad\text{(deployment hex of slot $j$)},\qquad
    e_j:=5(j-1)+1\quad\text{(hex of $E_j$)} ,
  \]
  the remaining three hexes of the block being empty. There are $k=n$ slots.
\item \textbf{Player army.} One creature type $C$ with $\mathit{att}=\mathit{def}=1$
  (that is $(\star)$), flat damage $d=1$, $\mathit{hp}=5$, $\mathit{spd}=2$,
  $\mathit{flags}=\emptyset$, $\mathit{value}=0$, and stock exactly $B$.
\item \textbf{Defence.} For each $j$, the stack $E_j$ is \textbf{one} creature of a type
  with $\mathit{att}=\mathit{def}=1$, flat damage 1, $\mathit{hp}=a_j$, $\mathit{spd}=1$,
  $\mathit{flags}=\emptyset$, $\mathit{value}=a_j$, at $e_j$, playing $(\ddagger)$.
\item \textbf{Question.} $R=1$, $W=B$.
\end{itemize}

\noindent Since every $a_i\ge 1$ and $n\ge 1$ we have $B\ge 1$, so the stock and
$W=B\in\mathbb{Z}_{>0}$ are legal. The board is listed hex by hex in $5n$ cells and the
numbers $a_j$ occur only as hit points and values, in binary, so $G(a)$ is computable in
time polynomial in the binary encoding of $a$. (A map polynomial in the binary length is
a fortiori polynomial in the unary one; Theorem 1 is \emph{weak} hardness not because the
map fails in unary but because Proposition 1.1's pseudo-polynomial algorithm decides the
unary problem.)

\paragraph{Totality.} The three-way branch of the conventions section, instantiated:
route to $G_{\mathrm{no}}$ every encoding that is malformed, or has some $a_i\le 0$, or has
$\sum_i a_i$ odd---in the last case no $S$ can sum to the non-integer $\sum_i a_i/2$, so a
no-certificate is at hand; route the empty encoding $n=0$, which is a \PARTITION{}
yes-instance because the empty subset sums to $0=B$, to the fixed instance
$G\bigl((1,1)\bigr)$, a yes-instance of \ALLOC{} by Lemma E.6; and otherwise output $G(a)$.

\begin{lemEsepone}[separation in the corridor]
In $G(a)$, $\mathrm{dist}(p_j,e_j)=1$ for every $j$, and
$\mathrm{dist}(p_j,e_{j'})=|5(j-j')-1|\ge 4$ for $j\ne j'$. A player stack has strike radius
$\mathit{spd}+1=3$ (R11), so in \emph{every} position of round 1 a stack standing on $p_j$
can strike no enemy other than $E_j$, and can strike $E_j$ whenever $E_j$ is alive.
\end{lemEsepone}

\noindent\emph{Proof.} On a one-row board R2 degenerates to $|\delta x|$, so
$\mathrm{dist}(p_j,e_{j'})=|5(j-1)-5(j'-1)-1|=|5(j-j')-1|$, which is 1 at $j'=j$ and, for
$|j-j'|\ge 1$, at least $|5-1|=4$. By Lemma E.3 a stack of speed 2 on $p_j$ can strike only
enemies at distance at most 3, and $E_j$ is the only one. Conversely $e_j$ is adjacent to
$p_j$, so while $E_j$ is alive \texttt{WALK\_AND\_ATTACK} against it with the empty walk and
approach hex $p_j$ is legal (R11), whatever else stands on the board. \qedbox

\smallskip
The two numbers here are the reduction's one tight margin, and the phrase ``consecutive
blocks are 4 apart'', which an earlier body draft used, named neither of them: consecutive blocks are \textbf{5}
apart (block $j$ starts at $5(j-1)$ and block $j+1$ at $5j$); what equals 4 is
$\min_{j'\ne j}\mathrm{dist}(p_j,e_{j'})$, attained at $j'=j-1$, the forward gap being 6.
Block width 5 is minimal: at width 4 the backward gap would be $|4-1|=3$, exactly the strike
radius, and slot $j$ would reach $E_{j-1}$.

\begin{lemEkillone}[the threshold at $E_j$]
Fix an allocation $(c_1,\dots,c_n)$ with $\sum_j c_j\le B$ and any play of round 1.
Let $A_j$ be the set of slots whose stack struck $E_j$ during the round. Then
\begin{enumerate}
\item $A_j\subseteq\{j\}$, and the sets $A_1,\dots,A_n$ are pairwise disjoint;
\item the total damage $E_j$ absorbs is at most $\sum_{i\in A_j}c_i$, so $E_j$ dead implies
  $\sum_{i\in A_j}c_i\ge a_j$;
\item if $c_j\ge a_j$ and the stack of slot $j$ issues \texttt{WALK\_AND\_ATTACK} against
  $E_j$ without waiting, $E_j$ dies.
\end{enumerate}
\end{lemEkillone}

\noindent\emph{Proof.} (1) By Lemma E.1 every blow a player stack delivers is issued at its
terminal action from its deployment hex, so by Lemma E.4 the stack of slot $i$ can only
strike $E_i$; hence $A_j\subseteq\{j\}$ and distinct $A_j$ are disjoint. (Disjointness in
fact needs no geometry: a stack delivers at most one blow, so it belongs to at most one
$A_j$.)

(2) By the one-round lemma of Section~\ref{sec:policy} every blow delivers at most its
nominal damage, and under $(\star)$ the nominal damage of the stack of slot $i$ is
$c_i\cdot d=c_i$ ($d=1$). Damage accumulates in $E_j$'s pool and excess is discarded (R8), so
the absorbed total is at most $\sum_{i\in A_j}c_i$. The stack $E_j$ is a single creature, so
by R5 its pool is $(\mathit{fullUnits},\mathit{firstHPleft})=(0,a_j)$ and
$\mathrm{count}=1$ (R6); by R8, $\mathrm{kills}(D)=0$ for $D<a_j$ and
$\mathrm{kills}(D)=\min(1+\lfloor (D-a_j)/a_j\rfloor,1)=1$ otherwise. So $E_j$ dies exactly
when the damage it absorbs reaches $a_j$, and death forces $\sum_{i\in A_j}c_i\ge a_j$.

(3) A non-waiting blow in round 1 lands in the \texttt{NORMAL} phase, strictly before every
enemy's postponed \texttt{DEFEND}, so it meets the un-raised defence
(Section~\ref{sec:policy}) and delivers its full nominal $c_j\ge a_j$ (here $c_j\ge a_j\ge 1$
and $d=1$, so $(\star)$ applies). By Lemma E.4 the approach hex may be $p_j$ itself, at
the approach reached by an empty walk, $p_j$ being adjacent to $E_j$ by Lemma E.4, so the strike is legal whatever else is on the board; and by Lemma E.2
no player creature dies, so the stack is alive to take its action (R9). \qedbox

\begin{lemEyesone}[\PARTITION{} yes $\Longrightarrow$ game yes]
If $a$ is a \PARTITION{} yes-instance then $G(a)$ is a yes-instance.
\end{lemEyesone}

\noindent\emph{Proof.} Let $S\subseteq[n]$ with $\sum_{j\in S}a_j=B$. Allocate $c_j:=a_j$
for $j\in S$ and $c_j:=0$ otherwise; the total is $B$, exactly the stock, so the allocation
is feasible, and each slot receives at most one type because there is only one. Let every
nonempty stack issue \texttt{WALK\_AND\_ATTACK} against the enemy of its own block, from its
own hex, without waiting. Each such $E_j$ is struck only by slot $j$, so it stands at full
pool when struck; by Lemma E.5(3) it dies. The destroyed value is
$\sum_{j\in S}a_j=B=W$. \qedbox

\begin{lemEnoone}[game yes $\Longrightarrow$ \PARTITION{} yes]
If $G(a)$ is a yes-instance then $a$ is a \PARTITION{} yes-instance.
\end{lemEnoone}

\noindent\emph{Proof.} Fix a feasible allocation $(c_1,\dots,c_n)$, $\sum_j c_j\le B$, and a
play destroying value at least $W=B$. The only creatures carrying value are the $E_j$, of
value $a_j$ each, so with $S:=\{j: E_j \text{ dead at the end of round 1}\}$ the destroyed
value is exactly $\sum_{j\in S}a_j\ge B$. By Lemma E.5(2), $\sum_{i\in A_j}c_i\ge a_j$ for
every $j\in S$, and by Lemma E.5(1) the $A_j$ are pairwise disjoint subsets of $[n]$. Hence
\begin{equation*}
  B\ \le\ \sum_{j\in S}a_j\ \le\ \sum_{j\in S}\ \sum_{i\in A_j}c_i\ \le\ \sum_{i=1}^{n}c_i
  \ \le\ B ,
  \tag{$\ddagger_1$}
\end{equation*}
so every inequality is an equality; in particular $\sum_{j\in S}a_j=B$ and $S$ solves
\PARTITION. \qedbox

\smallskip
\noindent\textbf{Two remarks on $(\ddagger_1)$.} First, the chain uses \emph{no reach
hypothesis}: only that the striker sets are disjoint (Lemma E.1) and that each dead $E_j$
absorbed at least $a_j$ from them. Lemma E.4 is therefore needed for the yes-direction
(where slot $j$ must be adjacent to $E_j$) but not for the no-direction. Consequently an
earlier version of Section~\ref{sec:generalized}, which warned that widening the player's
speed would break Theorem 1, was too strong: translation and speed-widening both preserve the
yes-direction, and the no-direction never looked at the board. What a wider speed breaks is
Proposition 1.1's hypothesis, which is a different statement about a different family---and
that is what Section~\ref{sec:generalized} now says.
Second, the chain bounds the \emph{total} damage $E_j$ absorbs, not the size of one blow;
an earlier body draft's ``its blow delivers at most $c_j$'' bounded the wrong quantity even though the
number is the same.

\smallskip
\noindent\textbf{Proof of Theorem 1.} Membership in NP is Lemma 2.1. The map
$a\mapsto G(a)$, extended by the totality branch above, is computable in time polynomial in
the binary encoding of $a$ and, by Lemmas E.6 and E.7 together with the two fixed instances
of the branch, sends yes-instances to yes-instances and no-instances to no-instances. The
constructed instances have $R=1$, a single player creature type, one creature per enemy
stack, no obstacles and one row. \PARTITION{} is NP-complete~\cite{Kar72}, so \ALLOC{} is
NP-complete on this family. \qedbox

\subsection*{Proposition 1.1, and what ``matching reach'' has to mean}

The family Theorem 1 constructs admits a pseudo-polynomial algorithm, which is what makes
the weak hardness of Theorem 1 tight. Stating the algorithm's hypothesis correctly takes
one definition, because the naive reading---``at the start, slot $j$ reaches exactly one
enemy''---is a statement about the starting position, while the algorithm needs a statement
about the whole round.

\begin{defEmatch}[persistent matching reach]
An instance with $R=1$, defence $(\ddagger)$ and $k$ slots has \textbf{persistent matching
reach} if there is a bijection $j\mapsto E_j$ from the slots onto the enemy stacks such that
for every feasible allocation, in every position of round 1 reachable from the starting
position that allocation induces by legal player actions
and the defence's $(\ddagger)$ actions, for every slot $j$ whose stack has \emph{not yet
taken its terminal action}, the set of enemies that stack can strike is exactly $\{E_j\}$.
\end{defEmatch}

A stack that has not taken its terminal action has not moved and has not struck (Lemma E.1),
so it still stands on $p_j$; ``the set of enemies it can strike'' is therefore the set of
enemies $E$ for which some free hex adjacent to $E$'s hex is within movement range of $p_j$
over the free hexes of the current position. The quantification over positions is the whole
point of the definition:

\begin{remark}
A dead unit stops blocking its hex (R4), so the free graph grows during the round and reach
is \emph{dynamic}. A hypothesis imposed only on the starting position would leave open that
a kill opens a doorway through which a second slot reaches a third slot's enemy---exactly the
leak that Theorem 3 closes with Lemma D.7. Definition E.8 closes it by fiat and Lemma E.9
discharges it, for the family at hand, with a criterion that never reads the free graph.
\end{remark}

\begin{lemEtest}[a metric test for Definition E.8]
Suppose the enemies are $E_1,\dots,E_k$ at hexes $e_1,\dots,e_k$, the slots have deployment
hexes $p_1,\dots,p_k$ and speeds $s_1,\dots,s_k$, and
\begin{enumerate}
\item $p_j$ is adjacent to $e_j$ for every $j$, and
\item $\mathrm{dist}(p_j,e_{j'})>s_j+1$ for all $j\ne j'$.
\end{enumerate}
Then the instance has persistent matching reach, with the bijection $j\mapsto E_j$.
\end{lemEtest}

\noindent\emph{Proof.} Fix a reachable position and a slot $j$ whose stack has not taken its
terminal action; it stands on $p_j$.

\emph{Containment.} If that stack strikes an enemy at hex $z$ then $\mathrm{dist}(p_j,z)\le
s_j+1$ by Lemma E.3, and hypothesis (2) rules out every $e_{j'}$ with $j'\ne j$. Note that
Lemma E.3 is insensitive to which hexes are free, so this holds in every reachable position,
however many hexes earlier kills have freed.

\emph{Non-emptiness.} By containment applied to every slot, no stack other than $j$'s can
ever strike $E_j$; and $j$'s stack has not struck. So $E_j$ is alive. Its hex is adjacent to
$p_j$ by (1), so \texttt{WALK\_AND\_ATTACK} against $E_j$ with the empty walk and approach
hex $p_j$ is legal (R11), and $E_j$ is strikable. \qedbox

\begin{lemEfam}[the family of Theorem 1 qualifies]
Every instance $G(a)$ built in the Theorem 1 section has persistent matching reach.
\end{lemEfam}

\noindent\emph{Proof.} Lemma E.4 gives $\mathrm{dist}(p_j,e_j)=1$ and
$\mathrm{dist}(p_j,e_{j'})\ge 4>3=\mathit{spd}+1$ for $j'\ne j$; apply Lemma E.9. \qedbox

\begin{propdp}
On the family of Theorem 1---a single creature type of flat damage $d\ge 1$ and stock $B$,
$R=1$, policy $(\ddagger)$, damage under $(\star)$, \textbf{one creature per enemy stack},
and \textbf{persistent matching reach} in the sense of Definition E.8---\ALLOC{} is decidable
in $O(k\,B)$ arithmetic operations and $O(B)$ working space after a polynomial-time
preprocessing pass (one breadth-first search per slot), hence in pseudo-polynomial time.
\end{propdp}

\noindent\emph{Proof.} Write $t_j$ for the hit points and $v_j$ for the value of the single
creature of $E_j$, and $b_j:=\lceil t_j/d\rceil$; $b_j$ is well defined because $d\ge 1$
(conventions section). We show that the optimum destroyed value equals
\begin{equation*}
  \mathrm{OPT}\ =\ \max\Bigl\{\ \sum_{j\in S}v_j\ :\ S\subseteq[k],\ \sum_{j\in S}b_j\le B\ \Bigr\},
  \tag{K}
\end{equation*}
a 0-1 knapsack over $k$ items, which the textbook dynamic program over (slot prefix,
consumed budget) solves by keeping one budget row: $O(kB)$ arithmetic operations on integers
of the input's bit length, and $O(B)$ stored values. That is the bound of the statement, and
it counts the dynamic program only: the thresholds $b_j$ and the bijection $j\mapsto E_j$ are
computed first, in a polynomial-time preprocessing pass that reads the board---one
breadth-first search per slot $j$, run in the starting position of the allocation that
deploys a single creature in slot $j$ and nothing else (feasible whenever $B\ge 1$; the case
$B=0$ is trivial). That starting position is a position reachable in the sense of
Definition E.8, so the set the search finds is exactly $\{E_j\}$. The whole algorithm is
therefore pseudo-polynomial---polynomial in the input length and in $B$---and not $O(kB)$
outright: an input whose board is large and whose $kB$ is small costs more than $kB$ to
read.

\emph{The optimum is at most $(\mathrm{K})$.} Fix a feasible allocation $(c_1,\dots,c_k)$,
$\sum_j c_j\le B$, and any play of round 1; let $S$ be the set of $j$ with $E_j$ dead. By
Lemma E.1 every blow is issued at its stack's terminal action from its deployment hex, so
Definition E.8 applies to it: the stack of slot $i$ strikes only $E_i$, hence $E_j$ is
struck only by slot $j$. By the one-round lemma every blow delivers \emph{at most} its
nominal damage, which is $c_j\,d$ under $(\star)$, and damage accumulates in the pool with
overkill discarded (R8). Since $E_j$ is a single creature of $t_j$ hit points, R5, R6 and R8
give $\mathrm{count}=1$, $\mathit{firstHPleft}=t_j$ and $\mathrm{kills}(D)=1$ exactly when
$D\ge t_j$. So $j\in S$ forces $c_j\,d\ge t_j$, i.e.\ $c_j\ge b_j$. Therefore
$\sum_{j\in S}b_j\le\sum_{j\in S}c_j\le B$ and the destroyed value $\sum_{j\in S}v_j$ is one
of the sums in $(\mathrm{K})$.

\emph{The optimum is at least $(\mathrm{K})$.} Let $S$ attain $(\mathrm{K})$. Allocate
$c_j:=b_j$ for $j\in S$ and $c_j:=0$ otherwise, feasible since $\sum_{j\in S}b_j\le B$.
Instruct every nonempty stack to issue \texttt{WALK\_AND\_ATTACK} against its own $E_j$,
without waiting; no stack waits, so every blow lands in the \texttt{NORMAL} phase and meets
the un-raised defence (Section~\ref{sec:policy}) and delivers its full nominal
$c_j\,d\ge t_j$ (here $c_j=b_j\ge 1$ and $d\ge 1$). Each such action is legal when it is
taken: the stack has not yet taken its terminal action, so by Definition E.8 the set of
enemies it can strike is exactly $\{E_j\}$---in particular $E_j$ is alive and an approach hex
is free, \emph{in that position}, whatever the earlier movers have done. This is where the
quantification over positions earns its keep: no separate argument is needed that an earlier
mover cannot block the route, because any position in which it did would be a reachable
position violating Definition E.8. Every $E_j$, $j\in S$, therefore dies and the destroyed
value is $\sum_{j\in S}v_j$. \qedbox

\smallskip
\noindent\textbf{Both extra hypotheses are load-bearing.} \emph{One creature per enemy
stack:} against a stack of several creatures a non-finishing blow still kills whole creatures
and still scores, the per-slot value is a staircase in $c_j$ rather than a threshold, and
$(\mathrm{K})$ is false---round 10 exhibited a matching-reach instance with six-creature
stacks on which the threshold rule returns 0 against a true optimum of 3 (open problem 4,
Section~\ref{sec:open}). \emph{Persistent matching reach:} it does not follow from the other
restrictions of Theorem 1's statement, not even on one row. Take hexes $0,\dots,7$; one
player type of flat damage 1, $\mathit{hp}=5$, $\mathit{spd}=1$, stock $B=2$; slot 1 at hex
1 and slot 2 at hex 7; single-creature enemies $E_1$ at hex 0 and $E_2$ at hex 2, each of 1
hit point and value 1; $R=1$, $W=2$. Every restriction in Theorem 1's statement holds. Slot
1 is adjacent to both enemies, so the reach is not a matching; slot 2 has strike radius 2
and reaches nothing. The true optimum is 1---slot 1 strikes once---while $(\mathrm{K})$ with
$b_1=b_2=1$ and $B=2$ returns 2, and the dynamic program, which indexes its items by slots
through the enemy each slot uniquely reaches, is not even well defined. What Proposition 1.1
shows tight is the family the reduction \emph{constructs}, not the family Theorem 1
quantifies over.

\subsection*{Theorem 2}

\begin{thmtwo}
\ALLOC{} is \textbf{strongly} NP-hard, already for $R=1$, no obstacles, no abilities, one
creature per enemy stack, and instances in which every player creature type has stock exactly
one.
\end{thmtwo}

\paragraph{Source problem.} \THREEPART{} (\cite[SP15]{GJ79}, strongly NP-complete): given
$3m$ positive integers $a_1,\dots,a_{3m}$ and a bound $T$ with $\sum_i a_i=mT$ and
$T/4<a_i<T/2$ for all $i$, is there a partition of $[3m]$ into $m$ triples each summing to
$T$? It remains NP-complete with all $a_i$ bounded by a polynomial in $m$, i.e.\ written in
unary.

\paragraph{The instance $G_{3P}(a,T)$.} Given such an instance with $m\ge 1$:

\begin{itemize}
\item \textbf{Battlefield.} Three rows ($y\in\{0,1,2\}$) and $8m+2$ columns; no obstacles.
  For $g=1,\dots,m$ put $X_g:=8(g-1)+1$ and place
  \[
    e_g:=(X_g,1),\qquad
    q_g^1:=(X_g-1,\,1),\quad q_g^2:=(X_g,\,0),\quad q_g^3:=(X_g,\,2).
  \]
  There are $k=3m$ slots, with deployment hexes
  $q_1^1,q_1^2,q_1^3,\dots,q_m^1,q_m^2,q_m^3$.
\item \textbf{Player army.} $3m$ creature types $C_1,\dots,C_{3m}$, with
  $\mathit{att}(C_i)=\mathit{def}(C_i)=1$ (that is $(\star)$), flat damage
  $\mathit{dmg}_{\min}=\mathit{dmg}_{\max}=a_i$, $\mathit{hp}=5$, $\mathit{spd}=2$,
  $\mathit{flags}=\emptyset$, $\mathit{value}=0$, and \textbf{stock exactly one}.
\item \textbf{Defence.} $E_g$ is \textbf{one} creature with $\mathit{att}=\mathit{def}=1$,
  flat damage 1, $\mathit{hp}=T$, $\mathit{spd}=1$, $\mathit{flags}=\emptyset$,
  $\mathit{value}=1$, at $e_g$, playing $(\ddagger)$.
\item \textbf{Question.} $R=1$, $W=m$.
\end{itemize}

\noindent Every number is bounded by a polynomial in $m$ once the $a_i$ are, and the board
has $3(8m+2)=O(m)$ hexes, so the construction is polynomial even under \textbf{unary}
encoding of the $a_i$---which is what makes the hardness strong. The board is wide enough:
$q_1^1$ sits in column $0$ and $e_m$ in column $X_m=8m-7\le 8m+1$.

\paragraph{Totality.} As in the conventions section: route to $G_{\mathrm{no}}$ every
encoding that is malformed, has a number of items not divisible by 3, has some $a_i\le 0$,
or fails $\sum_i a_i=mT$ or $T/4<a_i<T/2$---each check is a polynomial-time no-certificate
under the convention that an encoding outside the problem's domain is a no; route the empty
encoding $m=0$, a yes-instance, to the fixed instance $G_{3P}\bigl((1,1,1),3\bigr)$, a
yes-instance of \ALLOC{} by Lemma E.13; otherwise output $G_{3P}(a,T)$.

\begin{lemEseptwo}[three seats, and separation at distance 7]
In $G_{3P}(a,T)$ each of $q_g^1,q_g^2,q_g^3$ is a neighbour of $e_g$, the three are distinct,
and for $g'\ne g$
\[
  \mathrm{dist}(q_g^r,\,e_{g'})\ \ge\ 7 \qquad (r=1,2,3),
\]
with equality attained, e.g.\ at $q_{g+1}^1=(X_g+7,1)$ against $e_g$. Since a player stack
has strike radius $\mathit{spd}+1=3$ (R11), in \emph{every} position of round 1 a stack
standing on $q_g^r$ can strike no enemy other than $E_g$, and can strike $E_g$ whenever
$E_g$ is alive.
\end{lemEseptwo}

\noindent\emph{Proof.} Row 1 is odd, so by R1 the six neighbours of $(X_g,1)$ are
$(X_g\pm 1,1)$, $(X_g-1,0)$, $(X_g,0)$, $(X_g-1,2)$, $(X_g,2)$; the three hexes
$q_g^1,q_g^2,q_g^3$ are among them and are pairwise distinct.

For the separation, pass to the axial coordinates $A(x,y)=x+\lfloor y/2\rfloor$ of R2, in
which the group-$g$ hexes read
\begin{align*}
  A(e_g)&=X_g\ \ (y=1), & A(q_g^1)&=X_g-1\ \ (y=1),\\
  A(q_g^2)&=X_g\ \ (y=0), & A(q_g^3)&=X_g+1\ \ (y=2),
\end{align*}
and put $D:=X_{g'}-X_g=8(g'-g)$, so $|D|\ge 8$. Then $(\delta A,\delta y)$ from $q_g^r$ to
$e_{g'}$ is $(D+1,0)$, $(D,1)$, $(D-1,-1)$ for $r=1,2,3$, and R2 gives

\begin{center}
\begin{tabular}{@{}llll@{}}
\toprule
 & $(\delta A,\delta y)$ & $D\ge 8$ & $D\le -8$ \\
\midrule
$r=1$ & $(D+1,\,0)$  & $\max(D+1,0)=D+1\ge 9$ & $|D+1|+0=|D|-1\ge 7$ \\
$r=2$ & $(D,\,1)$    & $\max(D,1)=D\ge 8$     & $|D|+1\ge 9$ \\
$r=3$ & $(D-1,\,-1)$ & $(D-1)+1=D\ge 8$       & $\max(|D-1|,1)=|D|+1\ge 9$ \\
\bottomrule
\end{tabular}
\end{center}

\noindent(the mixed-sign rows use the $|\delta A|+|\delta y|$ branch, the others the
$\max$ branch). The minimum is 7, attained at $r=1$, $D=-8$, i.e.\ at $q_{g+1}^1$ against
$e_g$. Since $7>3$, Lemma E.3 leaves $E_g$ as the only enemy a stack on $q_g^r$ can ever
strike. \qedbox

\smallskip
An earlier body draft said ``at distance at least 6''---true but slack, and its derivation
subtracted a unit twice; the computation above is in the paper's own metric, is tight, and
is what the body now states. The bound is machine-checked on every constructed instance in the worst case for the player,
with all $3m$ slots occupied (\src{brute\_force.py}, \texttt{check\_geometry\_3partition}).

\begin{lemEacctwo}[damage accounting]\sloppy
Fix an allocation and any play of round 1 of $G_{3P}(a,T)$. For $g\in[m]$ let $S_g\subseteq[3m]$
be the set of types whose stack struck $E_g$. Then the $S_g$ are pairwise disjoint,
$|S_g|\le 3$, the damage $E_g$ absorbs is at most $\sum_{i\in S_g}a_i$, and $E_g$ is dead at
the end of the round only if $\sum_{i\in S_g}a_i\ge T$.
\end{lemEacctwo}

\noindent\emph{Proof.} A slot is homogeneous and every type has stock one, so a slot holds at
most one creature and a type occupies at most one slot. By Lemma E.1 each stack delivers at
most one blow, from its own deployment hex, and by Lemma E.11 that blow can only reach the
enemy of its own group; hence each type lies in at most one $S_g$ (disjointness) and only
the three slots of group $g$ feed $S_g$, so $|S_g|\le 3$. A stack of type $C_i$ holds one
creature, so under $(\star)$ its nominal blow is $1\cdot a_i=a_i$, and by the one-round lemma
of Section~\ref{sec:policy} it delivers at most that. (At most, not strictly less: the
formula clamps at 1 and $a_i=1$ is legal, and only the inequality is used.) No further damage
reaches $E_g$: under $(\ddagger)$ no enemy initiates, so no player stack ever delivers a
retaliation (Lemma E.1). Damage accumulates in $E_g$'s pool with overkill discarded (R8), and
$E_g$ is a single creature, so R5 and R6 give
$(\mathit{fullUnits},\mathit{firstHPleft})=(0,T)$ and $\mathrm{count}=1$, whence by R8
$\mathrm{kills}(D)=0$ for $D<T$ and $\mathrm{kills}(D)=\min(1+\lfloor(D-T)/T\rfloor,1)=1$
for $D\ge T$: death requires the absorbed total, hence $\sum_{i\in S_g}a_i$, to reach $T$.
\qedbox

\smallskip
Three independent facts sit inside that proof and each is used again below: a one-creature
stack of $C_i$ delivers at most $a_i$; each stack strikes at most once; and \emph{no player
stack ever delivers retaliation damage}. The third is the one whose absence produced the
paper's recorded first error (Section~\ref{sec:errors}).

\begin{lemEyestwo}[\THREEPART{} yes $\Longrightarrow$ game yes]
If $(a,T)$ is a \THREEPART{} yes-instance then $G_{3P}(a,T)$ is a yes-instance.
\end{lemEyestwo}

\noindent\emph{Proof.} Let $\{G_1,\dots,G_m\}$ be a partition of $[3m]$ into triples with
$\sum_{i\in G_g}a_i=T$. Allocate the three types of $G_g$ to the three slots
$q_g^1,q_g^2,q_g^3$, one creature each---feasible, since each type has stock one and each
slot receives one type. Instruct every stack to issue \texttt{WALK\_AND\_ATTACK} against the
enemy of its own group, from its own hex, without waiting. We check that this play is legal
and that it kills every $E_g$; each item is a step the body's sketch left implicit.

\emph{Every stack acts.} Under $(\ddagger)$ no enemy initiates, and by Lemma E.2 the only
damage a player stack takes is one retaliation of 1 point against $\mathit{hp}=5$, delivered
strictly after that stack's own blow. So no player stack is killed or weakened before acting,
and every living stack takes its terminal action once in the round (R9).

\emph{Every blow is legal.} $q_g^r$ is adjacent to $e_g$ (Lemma E.11), so the approach hex is
the stack's own hex and the walk is empty (R11); no ally can block it and no earlier mover can
occupy it, since no stack in this play ever leaves its deployment hex.

\emph{Every blow delivers exactly $a_i$.} No player stack waits, so all $3m$ player blows land
in the \texttt{NORMAL} phase, while under $(\ddagger)$ each $E_g$ spends its \texttt{NORMAL}
activation on \texttt{WAIT} and issues its terminal \texttt{DEFEND} only in the later
\texttt{WAIT} phase (R9, Section~\ref{sec:policy}). Every blow therefore meets the un-raised
defence $\mathit{def}(E_g)=1$, so $\Delta=1-1=0$, $f_{\mathrm{att}}=f_{\mathrm{def}}=1$, and
by $(\star)$ the blow of $C_i$ delivers $\max(1,\lfloor 1\cdot a_i\rfloor)=a_i$ (R7). (This
holds regardless of relative speeds; that the player is faster is not needed.)

\emph{Every enemy is alive when struck, and dies.} The blows aimed at $E_g$ are the three
$a_i$, $i\in G_g$, and $a_i<T/2$, so any two of them sum to less than $T$: $E_g$ survives its
first two blows and the third striker finds it alive. Damage accumulates in the pool (R8), and
after the third blow the accumulated total is $\sum_{i\in G_g}a_i=T$, so
$\mathrm{kills}=1$ by Lemma E.12's computation and $E_g$ dies.

All $m$ enemies die and the destroyed value is $m\cdot 1=W$. \qedbox

\begin{lemEnotwo}[game yes $\Longrightarrow$ \THREEPART{} yes]
If $G_{3P}(a,T)$ is a yes-instance then $(a,T)$ is a \THREEPART{} yes-instance.
\end{lemEnotwo}

\noindent\emph{Proof.} Each enemy creature has value 1 and there are $m$ of them, so a
destroyed value of at least $W=m$ forces all $m$ to die. By Lemma E.12,
$\sum_{i\in S_g}a_i\ge T$ for every $g$, with the $S_g$ pairwise disjoint subsets of $[3m]$.
Summing,
\[
  mT\ \le\ \sum_{g=1}^{m}\ \sum_{i\in S_g}a_i\ \le\ \sum_{i=1}^{3m}a_i\ =\ mT ,
\]
so every inequality is tight: $\sum_{i\in S_g}a_i=T$ for each $g$, and the $S_g$ cover
$[3m]$. Finally $T/4<a_i<T/2$ forces $|S_g|=3$---two elements sum to less than $T$, four to
more---so $\{S_1,\dots,S_m\}$ is a 3-partition. (The $S_g$ are allowed to be proper subsets
of a group's three types, since slots may be left empty; the counting is what rules that
out.) \qedbox

\smallskip
\noindent\textbf{Proof of Theorem 2.} Lemmas E.13 and E.14 give the equivalence, and the
totality branch extends it to every encoding. The construction is computable in time
polynomial in the unary encoding of $(a,T)$ and every number in the output is bounded by a
polynomial in $m$ and $\max_i a_i$, so it is a pseudo-polynomial---indeed polynomial under
unary encoding---reduction; since \THREEPART{} is strongly NP-complete~\cite{GJ79}, \ALLOC{}
is strongly NP-hard on this family. The instances have $R=1$, no obstacles, no abilities, one
creature per enemy stack, and every player type of stock one. \qedbox

\begin{remark}
\textbf{The three rows are not decoration.} In one row a hex has two neighbours, so a third
stack could not reach $E_g$ without walking through a hex occupied by an ally, and occupied
hexes are not enterable (R4). The bug was found by the geometry self-check that verifies
Lemma E.11 on the \emph{built} instance with all slots occupied (the geometry self-check of
the Theorem 1--2 suite, Section~\ref{sec:layers}).
\end{remark}

\subsection*{Theorem 4 and its corollaries}

\begin{thmfour}
\ALLOC{} is \textbf{strongly} NP-hard already for instances with $R=1$; \textbf{no obstacles
and no abilities of any kind}; every player creature type of stock one, one creature per enemy
stack; a rectangular open battlefield of six rows and $4m+2$ columns; and \textbf{complete
reachability}---in the starting position, with every slot occupied, every player stack can
attack every enemy stack.
\end{thmfour}

\begin{corfourone}
The same instances are hard with the allocation \emph{given}. \PLAY{} is strongly NP-hard on
obstacle-free boards with complete reachability.
\end{corfourone}

\begin{corfourtwo}[hit-point objective]
On the same instances, replace the objective by \textbf{total enemy hit points removed}---the
damage the enemy stacks absorb, with overkill discarded---and the target by
$W_{\mathrm{hp}}=mT$. The problem remains strongly NP-hard, with the allocation free or given.
\end{corfourtwo}

\paragraph{The instance $G_F(a,T)$.} From a \THREEPART{} instance $(a_1,\dots,a_{3m};T)$ with
$m\ge 1$:

\begin{itemize}
\item \textbf{Battlefield.} $h=6$ rows ($y\in\{0,\dots,5\}$) and $w=4m+2$ columns;
  \textbf{no obstacles}.
\item \textbf{Defence.} For $g=1,\dots,m$ put $X_g:=4g-2$ and place $E_g$, \textbf{one}
  creature, at $e_g:=(X_g,3)$, of a type with $\mathit{att}=\mathit{def}=1$, flat damage 1,
  $\mathit{hp}=T$, $\mathit{spd}=1$, $\mathit{flags}=\emptyset$, $\mathit{value}=1$, playing
  $(\ddagger)$.
\item \textbf{Player army.} $3m$ creature types $C_1,\dots,C_{3m}$ with
  $\mathit{att}(C_i)=\mathit{def}(C_i)=1$, flat damage
  $\mathit{dmg}_{\min}=\mathit{dmg}_{\max}=a_i$, $\mathit{hp}=5$,
  $\mathit{spd}=s:=w+h=4m+8$, $\mathit{flags}=\emptyset$, $\mathit{value}=0$, and stock
  exactly one. There are $k=3m$ slots along the top row, $p_j:=(j-1,0)$ for $j=1,\dots,3m$.
\item \textbf{Question.} $R=1$, $W=m$.
\end{itemize}

\noindent The board fits: $X_1=2$ and $X_m=4m-2\le w-3$, so every $e_g$ has all six
neighbours on the board; and $3m-1\le w-1$, so the deployment row is long enough. The board
has $6(4m+2)=O(m)$ hexes and every number is polynomial in $m$ once the $a_i$ are, so the
construction is polynomial under unary encoding. The totality branch is that of the Theorem 2
section verbatim, with $G_F\bigl((1,1,1),3\bigr)$ as the fixed yes-instance. One
qualification to the statement: its family clause---six rows, $4m+2$ columns, complete
reachability---describes the image of every well-formed encoding; the degenerate encodings
map to the fixed no-instance $G_{\mathrm{no}}$ of Lemma D.4, a $1\times 1$ board, which is
not of that shape.

Three named hexes per enemy carry all the geometry:
\begin{equation*}
  q_g^1:=(X_g-1,\,2),\qquad q_g^2:=(X_g,\,2),\qquad q_g^3:=(X_g+1,\,3).
  \tag{Q}
\end{equation*}

\begin{lemEseatsfour}[the three approach hexes]
Row 3 is odd, so by R1 the six neighbours of $e_g=(X_g,3)$ are
\begin{equation*}
  (X_g-1,3),\ (X_g+1,3),\ (X_g-1,2),\ (X_g,2),\ (X_g-1,4),\ (X_g,4).
  \tag{N}
\end{equation*}
The three hexes of $(\mathrm{Q})$ are among them, and the $3m$ hexes
$\{q_g^r: g\in[m],\,r\in\{1,2,3\}\}$ are pairwise distinct and disjoint from the deployment
row.
\end{lemEseatsfour}

\noindent\emph{Proof.} With $\epsilon=1$ for the odd row 3, R1 gives upper neighbours at
columns $x-1$ and $x$ of row 2 and lower neighbours at columns $x-1$ and $x$ of row 4, plus
the two same-row hexes; that is $(\mathrm{N})$, and $q_g^1,q_g^2,q_g^3$ appear in it. Within a
group the three hexes differ in column or row. Across groups $|X_g-X_{g'}|\ge 4$, while the
columns used by group $g$ lie in $\{X_g-1,X_g,X_g+1\}$, so no two groups share a hex. All
$q_g^r$ lie in rows 2 and 3, the $p_j$ in row 0. \qedbox

\smallskip
We used the offset convention of R1 in the form stated in the conventions section: from an
\emph{odd} row the two upper neighbours sit at columns $x-1$ and $x$; from an \emph{even} row
at columns $x$ and $x+1$. Every neighbour computation below is in that convention.

\smallskip
Call a play of round 1 \textbf{attack-only} if every player stack, on its turn, either passes
or issues \texttt{WALK\_AND\_ATTACK}---no \texttt{MOVE}-only action, no \texttt{WAIT}, no
\texttt{DEFEND}. Lemmas E.16--E.19 are used \emph{only} in the yes-direction, where the play
is ours to choose, so restricting them to attack-only plays costs nothing; the no-direction
(Lemma E.22) is geometry-free and quantifies over every play of the full model.

\begin{lemErow}[row 1 stays clear]
In any attack-only play of round 1 of $G_F$, every occupied hex lies in row 0, or is an enemy
hex, or is adjacent to an enemy hex. In particular \textbf{rows 1 and 5 are free throughout}.
\end{lemErow}

\noindent\emph{Proof.} Enemy stacks sit on $e_g$ in row 3 and never move: under $(\ddagger)$
they issue \texttt{WAIT} and then \texttt{DEFEND} in place, and neither changes a hex
(Section~\ref{sec:policy}). A player stack starts on $p_j$ in row 0, and in an attack-only
play the only action that changes its hex is \texttt{WALK\_AND\_ATTACK}, whose destination is
adjacent to the struck enemy (R11), hence adjacent to some $e_g$, hence by $(\mathrm{N})$ in
rows 2, 3 or 4. Each stack takes one terminal action (R9), so no other hex is ever entered.
Rows 1 and 5 contain no $e_g$ and, by $(\mathrm{N})$, no neighbour of any $e_g$. \qedbox

\begin{lemEroute}[routing]
Fix $g$ and $r\in\{1,2,3\}$. In any position of an attack-only play of round 1 in which
$E_g$ is alive and $q_g^r$ is free, a player stack still standing on its deployment hex $p_j$ can move to
$q_g^r$ and strike $E_g$. The walk has length at most $w+2=4m+4<s$, and this holds
independently of the order in which the stacks are activated.
\end{lemEroute}

\noindent\emph{Proof.} By Lemma E.16 every hex of row 1 is free. Row 0 is even, so by R1 the
two lower neighbours of $p_j=(j-1,0)$ are $(j-1,1)$ and $(j,1)$: the stack leaves row 0 in one
step and may then walk along row 1 freely, at most $w-1$ steps. Row 1 is odd, so the two lower
neighbours of $(x,1)$ are $(x-1,2)$ and $(x,2)$; row 2 is even, so the two lower neighbours of
$(x,2)$ are $(x,3)$ and $(x+1,3)$. Hence
\begin{itemize}
\item $q_g^1=(X_g-1,2)$ is entered from $(X_g-1,1)$, one step down from row 1;
\item $q_g^2=(X_g,2)$ is entered from $(X_g,1)$, likewise;
\item $q_g^3=(X_g+1,3)$ is entered from $(X_g+1,2)$, which is entered from $(X_g+1,1)$. The
  intermediate hex $(X_g+1,2)$ is \emph{always} free: its neighbours are, by R1 for the even
  row 2, $(X_g,2)$, $(X_g+2,2)$, $(X_g+1,1)$, $(X_g+2,1)$, $(X_g+1,3)$ and $(X_g+2,3)$, and no
  enemy hex $(X_{g'},3)$ is among them, since $X_{g'}-X_g\in 4\mathbb{Z}$ excludes
  $X_{g'}\in\{X_g+1,X_g+2\}$; so by Lemma E.16 it is never occupied.
\end{itemize}
Each route is one step out of row 0, at most $w-1$ steps along row 1, and at most two steps
down, so at most $w+2=4m+4$ steps, all over free hexes, well within $s=4m+8$ (R10). By R11 the
stack may therefore move to $q_g^r$ and strike the adjacent $E_g$ (Lemma E.15). Nothing in the
route depends on which other stacks have already acted, beyond the freeness of $q_g^r$, which
is the hypothesis. \qedbox

\begin{lemEreachfour}[complete reachability]
In the starting position of $G_F$, with all $3m$ slots occupied, every player stack can attack
every enemy stack.
\end{lemEreachfour}

\noindent\emph{Proof.} At the start the occupied hexes are exactly row 0 and the $e_g$, so
every $q_g^1$ is free; apply Lemma E.17. \qedbox

\begin{lemEreal}[simultaneous realizability]
Let $\varphi:[3m]\to[m]$ satisfy $|\varphi^{-1}(g)|=3$ for every $g$. Then there is an
attack-only play of round 1 of $G_F$, with all $3m$ slots occupied, in which every stack $j$
strikes $E_{\varphi(j)}$.
\end{lemEreal}

\noindent\emph{Proof.} For each $g$ let $j_g^1<j_g^2<j_g^3$ be the three slots with
$\varphi(j)=g$, and instruct stack $j_g^r$ to move to $q_g^r$ and strike $E_g$. This assigns
the $3m$ stacks to the $3m$ pairwise distinct hexes $q_g^r$ bijectively (Lemma E.15), so no
two stacks are sent to the same hex, and by construction the only stack that ever enters
$q_g^r$ is $j_g^r$. Hence $q_g^r$ is free when $j_g^r$ acts, and Lemma E.17 makes the move and
the strike legal; both are independent of the activation order, so the play is well defined
whatever order R9 imposes.

Two further points make the play realize $\varphi$ rather than merely attempt it. Every stack
gets its action: under $(\ddagger)$ no enemy initiates, and by Lemma E.2 a stack takes at most
one retaliation, of 1 point against $\mathit{hp}=5$, strictly after its own blow, so no stack
is killed before acting (R9). And every $E_g$ is alive when each of its three strikers acts:
the blows aimed at it are the three $a_i$ with $\varphi(i)=g$, and $a_i<T/2$, so any two of
them sum to less than $T=\mathit{hp}(E_g)$ and $E_g$ survives the first two. Note that this
argument uses only the \THREEPART{} promise $a_i<T/2$, so the lemma holds for \emph{every}
three-per-enemy $\varphi$, not only for those a 3-partition produces. \qedbox

\begin{lemEaccfour}[damage accounting]
Fix any allocation and any play of round 1 of $G_F$---not necessarily attack-only. For
$g\in[m]$ let $S_g\subseteq[3m]$ be the set of types whose stack struck $E_g$. Then the $S_g$
are pairwise disjoint, the nominal damage delivered to $E_g$ is at most $\sum_{i\in S_g}a_i$,
the damage $E_g$ absorbs is at most $\min\bigl(T,\ \sum_{i\in S_g}a_i\bigr)$---with equality
when no striker of $E_g$ waited; a waiting blow meets the postponed \src{DEFEND} bonus and
delivers at most nominal---under $(\star)$ the bonus is 1 and the factor $0.975$, so equality
holds exactly when $a_i=1$, where the clamp at 1 binds, and the blow is strictly less
otherwise---and $E_g$ is dead at the
end of the round only if $\sum_{i\in S_g}a_i\ge T$. If moreover no striker of $E_g$ waited and
$\sum_{i\in S_g}a_i\ge T$, then $E_g$ is dead.
\end{lemEaccfour}

\noindent\emph{Proof.} Each type has stock one and a slot is homogeneous, so each type
occupies at most one slot and each stack holds at most one creature. By Lemma E.1 each stack
delivers at most one blow in the round---one terminal action (R9, R11), and no retaliation
blow, because under $(\ddagger)$ no enemy initiates---so each stack strikes at most one enemy
and the $S_g$ are pairwise disjoint. By Lemma E.2 a player stack never loses a creature, so
the stack of $C_i$ has $\mathrm{count}=1$ when it strikes and its nominal blow is $a_i$ under
$(\star)$; by the one-round lemma of Section~\ref{sec:policy} it delivers at most that, and
exactly that if it did not wait. Damage accumulates in $E_g$'s pool and the excess is
discarded (R8), and $E_g$ is one creature of $T$ hit points, so by R5, R6 and R8 the absorbed
amount is $\min(T,\text{delivered})$ and $E_g$ dies exactly when the absorbed amount reaches
$T$. Both implications follow. \qedbox

\smallskip
Lemma E.20 is where the ``each type has stock one, so each stack strikes at most one
enemy'' is discharged properly. Stock one bounds the number of \emph{slots} a type occupies;
what bounds the blows is one terminal action per round together with the absence of enemy
initiative. As printed there, that justification would survive an attacking defence, which the
theorem does not.

\begin{lemEyesfour}[\THREEPART{} yes $\Longrightarrow$ game yes]
If $(a,T)$ is a \THREEPART{} yes-instance then $G_F(a,T)$ is a yes-instance, and the witness
uses the allocation $C_i\mapsto\text{slot } i$.
\end{lemEyesfour}

\noindent\emph{Proof.} Let $\{G_1,\dots,G_m\}$ be triples with $\sum_{i\in G_g}a_i=T$. Deploy
$C_i$ in slot $i$ (any injection would do; each type has stock one, so this is feasible) and
set $\varphi(i):=g$ for $i\in G_g$. By Lemma E.19 there is an attack-only play in which every
stack strikes its assigned enemy; no stack waits, so by Lemma E.20 each $E_g$ is delivered
exactly $\sum_{i\in G_g}a_i=T$ and dies. The destroyed value is $m\cdot 1=W$. \qedbox

\begin{lemEnofour}[game yes $\Longrightarrow$ \THREEPART{} yes]
If $G_F(a,T)$ is a yes-instance then $(a,T)$ is a \THREEPART{} yes-instance.
\end{lemEnofour}

\noindent\emph{Proof.} Enemy creatures have value 1 and there are $m$ of them, so a destroyed
value of at least $W=m$ forces all $m$ to die. By Lemma E.20, $\sum_{i\in S_g}a_i\ge T$ for
every $g$, with the $S_g$ pairwise disjoint subsets of $[3m]$; summing,
\[
  mT\ \le\ \sum_{g=1}^{m}\ \sum_{i\in S_g}a_i\ \le\ \sum_{i=1}^{3m}a_i\ =\ mT ,
\]
so every inequality is tight, the $S_g$ cover $[3m]$ and each sums to $T$; $T/4<a_i<T/2$ then
forces $|S_g|=3$. This direction uses no geometry: the board can only ever \emph{restrict}
which $S_g$ are achievable, so no freedom of movement can help. \qedbox

\smallskip
\noindent\textbf{Proof of Theorem 4.} Lemmas E.21 and E.22 give the equivalence and the
totality branch extends it to every encoding; the construction is polynomial under unary
encoding, and \THREEPART{} is strongly NP-complete~\cite{GJ79}, so \ALLOC{} is strongly
NP-hard on this family. The instances have $R=1$, no obstacles, no abilities, every player
type of stock one, one creature per enemy stack, a $6\times(4m+2)$ open rectangle, and
complete reachability by Lemma E.18. Membership in NP is Lemma 2.1, so the problem is strongly
NP-complete on this family. \qedbox

\smallskip
\noindent\textbf{Proof of Corollary 4.1.} Fix the allocation $C_i\mapsto\text{slot } i$.
Lemma E.21 uses exactly that allocation and Lemma E.22 never mentions the allocation, so the
equivalence survives with the allocation given as part of the input. \qedbox

\smallskip
\noindent\textbf{Proof of Corollary 4.2.} The objective is now the total number of
\textbf{enemy} hit points removed---the sum over enemy stacks of the damage they absorb, with
overkill discarded (R8). (Player hit points are removed too, by retaliation; they are not
counted, exactly as the Problem of Section~\ref{sec:problem} counts enemy values only.)

Fix any allocation and any play; let $A_g$ be the damage $E_g$ absorbs and $D_g$ the nominal
damage delivered to it---a waiting striker can deliver less than nominal, which is why the
two are named apart. By Lemma E.20 the $S_g$ are pairwise disjoint and
$A_g\le\min(T,D_g)$ with $D_g\le\sum_{i\in S_g}a_i$. Suppose the total absorbed reaches
$W_{\mathrm{hp}}=mT$. Each of the $m$ terms
$A_g$ is at most $T$, so all $m$ are exactly $T$; hence
$\sum_{i\in S_g}a_i\ge D_g\ge A_g=T$, and summing over the disjoint $S_g$ against
$\sum_i a_i=mT$ makes every inequality tight, so $\sum_{i\in S_g}a_i=T$ for every $g$ and the
$S_g$ cover $[3m]$; $T/4<a_i<T/2$ forces $|S_g|=3$ and $\{S_1,\dots,S_m\}$ is a 3-partition.

Conversely, the witness play of Lemma E.21 delivers exactly $T$ to each enemy, which each
absorbs in full, for a total of $mT$. The argument never mentions the allocation, so it holds
with the allocation free or given. \qedbox

\smallskip
\noindent\textbf{Where the hardness lives.} In $G_F$ every type has stock one and, by
Lemma E.18, every slot reaches every enemy, so every injection of types into slots is
equivalent: the multiset of blows available does not depend on which slot holds which type,
and by Lemma E.19 every three-per-enemy targeting is realizable from any of them. The decision
that encodes the 3-partition is the choice of targets---which is Corollary 4.1, and which is
what makes Theorem 4 the targeting-driven endpoint of the two-source claim.

\begin{remark}
The board does impose one thing: an enemy has six neighbours, so at most six stacks can strike
it in one round. Under $(\star)$ as completed in the conventions section this bound is
correct---no striker is killed by the retaliation it draws (Lemma E.2), so no seat is ever
vacated mid-round. It never binds here, three stacks per enemy, but it is why ``featureless''
means \emph{complete reachability plus local seat capacity} and not ``positions do not
exist''.
\end{remark}

\end{document}
